% Pass: Opus 5 (claude-opus-5), 2026-09-12 - sheets 1-12, first pass,
% unchecked against the photographs by a human.
% Pass: Opus 5 (claude-opus-5), 2026-09-12 - sheets 13-24, first pass,
% unchecked against the photographs by a human.
%
% The black notebook: a small two-ring leather notebook, 68 photographed
% sheets, of which the first sitting reads the first twelve. It is not a diary.
% It is a household and studio memorandum book, and its subject changes every
% page or two: the address and telephone at 3 Washington Square North, a
% phonograph and its spare parts, French phrases copied out with their English
% beside them, a fountain pen, the sizes of frames and of water colour paper,
% a consignment of her own water colours to a dealer, the recipients of a
% monograph, the insurance on the Truro cottage, which Broadway theatres have a
% usable second balcony, Edward Hopper's suit measurements, and a chronicle of
% where the two of them went in every year from 1961 back to 1936.
%
% Only one passage in the twelve sheets opens with a date, so there is one
% \entry{} - the monograph list of « Spring 1946 » on page 15. Everything else
% is given line for line, as Battle of Wash Sq. is. Several blocks carry a date
% at their foot rather than their head (« Jan. 20. 1954, » on page 1,
% « Jan 16, 60. » on page 2); those stand where they are written and head
% nothing.
%
% The year chronicle begins on page 19 and runs backwards - 1961 to 1957 on
% page 19, 1956 to 1946 on page 20, 1945 to 1936 on page 21 - and is given a
% \entry{} per year, because each year line heads its own memorandum. The
% leaves are not reordered, and the years therefore descend as the reader
% scrolls forward. Page 19's last line and page 20's first both record the same
% stay at Pacific Palisades in the winter of 1956-57; the repetition is hers.
%
% Four of the twelve sheets carry writing on one side only. Pages 7, 12 and 14
% are blank, and so is the front cover, which is plain black cloth with nothing
% written or printed on it; none of the four appears below.
%
% Three sheets are the hard ones. Page 3 has a block of five ink arrows and
% five strokes of red crayon whose lines were afterwards rubbed out, and
% nothing in them can be recovered. Page 15 is a list of some thirty names in
% two columns, the right-hand one written later, smaller, and across the left,
% and a good many of its names will not settle. Page 20 is the densest of the
% twelve, with an interlineation above almost every line.
%
% The second sitting reads the openings 22-23 to 44-45. The backwards year
% chronicle begun on page 19 runs on to the foot of page 24, where it reaches
% 1906 and stops; from page 25 the book returns to memoranda, and the subjects
% of this sitting are a chrysoprase and a coral necklace, tannic acid for
% burns, Who's Who in American Art, an eye test of 1952, three lists of
% photographs of her own pictures with their negative numbers, the staff of a
% dozen museums, the mailing lists for the three numbers of « Reality », the
% books Katherine Proctor gave her, two electric blankets, what was pulled down
% round Washington Square between 1953 and 1955, the sizes of her water colours,
% her own Who's Who entry copied out and boxed in red crayon, and a second
% chronicle - this one running forwards - of operations, hospitals and trips
% from 1941 to 1952.
%
% Pages 27 and 32 are blank and do not appear; page 26 carries the eye test and
% nothing faces it, and page 36 carries three lines. The openings are given
% under the Whitney's own descriptors, so a section heading names a blank page
% where the photograph does.
%
% Two sheets carry the sitting's only ruled tables. Page 26's eye test is ruled
% in three vertical strokes and is a \texttt{ledgertable}; page 34's two columns
% of names are divided by a drawn rule and are another. Page 28's « Arts +
% Letters » block stands in three columns that are not ruled, and is given as
% written, line by line, with a note - its middle column's dittos are read
% downward.
%
% Page 35 is the sheet to budget for, as page 15 was in the first sitting: some
% ninety names in two sittings of ink, many of them abbreviated to a surname and
% a tick, and a good many that will not settle.

% The Whitney gives this notebook no object number, so no \objectnumber line
% stands in the preamble.

\documentclass[11pt,a4paper]{article}
% The preamble shared by every transcription of the Hopper artist's ledgers.
%
% It rests on one idea, taken whole from the Grothendieck workbench this
% method comes from: **uncertainty compounds, it does not resolve.** A
% transcription that smooths over a doubtful figure has destroyed the only
% thing it was made for — a reader must be able to tell, at every point, what
% was on the leaf and what was supplied.
%
% A ledger sharpens that. In a manuscript of mathematics an invented word is a
% wrong word. Here an invented figure is a *sale that did not happen*, at a
% price nobody paid, to a buyer who never bought — and it will be cited,
% because a table looks like data in a way that prose does not.
%
% The same commands are recognised by `scripts/render.mjs`, which produces the
% site's reading view. A macro added here must be added there too, or rendering
% fails loudly, which is the wanted behaviour: a silently wrong rendering is
% worse than none.

% \ProvidesFile, not \ProvidesPackage: the transcriptions pull this in with
% % The preamble shared by every transcription of the Hopper artist's ledgers.
%
% It rests on one idea, taken whole from the Grothendieck workbench this
% method comes from: **uncertainty compounds, it does not resolve.** A
% transcription that smooths over a doubtful figure has destroyed the only
% thing it was made for — a reader must be able to tell, at every point, what
% was on the leaf and what was supplied.
%
% A ledger sharpens that. In a manuscript of mathematics an invented word is a
% wrong word. Here an invented figure is a *sale that did not happen*, at a
% price nobody paid, to a buyer who never bought — and it will be cited,
% because a table looks like data in a way that prose does not.
%
% The same commands are recognised by `scripts/render.mjs`, which produces the
% site's reading view. A macro added here must be added there too, or rendering
% fails loudly, which is the wanted behaviour: a silently wrong rendering is
% worse than none.

% \ProvidesFile, not \ProvidesPackage: the transcriptions pull this in with
% % The preamble shared by every transcription of the Hopper artist's ledgers.
%
% It rests on one idea, taken whole from the Grothendieck workbench this
% method comes from: **uncertainty compounds, it does not resolve.** A
% transcription that smooths over a doubtful figure has destroyed the only
% thing it was made for — a reader must be able to tell, at every point, what
% was on the leaf and what was supplied.
%
% A ledger sharpens that. In a manuscript of mathematics an invented word is a
% wrong word. Here an invented figure is a *sale that did not happen*, at a
% price nobody paid, to a buyer who never bought — and it will be cited,
% because a table looks like data in a way that prose does not.
%
% The same commands are recognised by `scripts/render.mjs`, which produces the
% site's reading view. A macro added here must be added there too, or rendering
% fails loudly, which is the wanted behaviour: a silently wrong rendering is
% worse than none.

% \ProvidesFile, not \ProvidesPackage: the transcriptions pull this in with
% \input{../preamble/hopper} rather than \usepackage, and \ProvidesPackage in
% an \input-ed file makes LaTeX warn on every compile that it "requested
% package `'" and got `hopper'. A warning nobody can act on is a warning
% everybody learns to scroll past, which is how the real ones get missed.
\ProvidesFile{hopper.sty}[2026/09/05 Transcriptions of the Hopper artist's ledgers]

\usepackage{array}
% longtable, not tabularx: a leaf of thirty-one exhibition lines must break
% across pages, and tabularx cannot be wrapped in an environment of our own —
% it scans the source for a literal \end{tabularx} and dies on \end{ledgertable}
% with « File ended while scanning use of \TX@get@body ». longtable has no such
% scan, and the wrapping column type below does the job tabularx's X would.
\usepackage{longtable}

% A wrapping column, `Y{0.3}` being three tenths of the measure.
%
% Prose columns must be declared this way and not as `l`. A table of `l`
% columns is set to its natural width, which can be wider than the page, and
% TeX issues **no warning at all** because nothing ever asked the row to fit:
% the first compile of Book I produced a page whose right-hand column ran off
% the paper with a clean log. A fixed-width column wraps, and if it still
% cannot fit, TeX says so — which is what `npm run pdf` then refuses to build
% over.
%
% Leave about 0.03 of the measure per column for the inter-column space: five
% columns can share 0.85, six can share 0.80. Getting it wrong is not a
% guessing game — `npm run pdf` reports the overfull box and how many points
% too wide it is.
\newcolumntype{Y}[1]{>{\raggedright\arraybackslash}p{#1\linewidth}}
\usepackage[english]{babel}
\usepackage{fontspec}
\usepackage{geometry}
\geometry{a4paper,margin=25mm,marginparwidth=28mm}
\usepackage{marginnote}
\usepackage{xcolor}
\usepackage{tikz}
% ulem, not soul. `soul`'s \st and \ul are built on a character-by-character
% scan that XeLaTeX's font machinery breaks: the document fails with an
% undefined control sequence rather than degrading. `normalem` stops it
% hijacking \emph, which the transcriptions use for underlining of Jo Hopper's
% own.
\usepackage[normalem]{ulem}
% ulem sets each underlined word in a box TeX can neither hyphenate nor break
% after, so a paragraph dense in \uncertain{} readings can reach a state where
% no legal breakpoint exists. \emergencystretch lets TeX loosen it on a third
% pass instead of overrunning: an overfull line in a transcription hides
% content.
\setlength{\emergencystretch}{3em}
\usepackage{hyperref}
\hypersetup{colorlinks=true,linkcolor=black,urlcolor=black,pdfborder={0 0 0}}

% --- Batch metadata ---------------------------------------------------------
% Taken as they stand by the rendering script, which shows them at the head of
% the reading view. They fix what the file claims to cover: without them a
% transcription floats unmoored in an archive of five hundred sheets.

\newcommand{\ledger}[1]{\gdef\hopLedger{#1}}
\newcommand{\ledgertitle}[1]{\gdef\hopLedgerTitle{#1}}
% A notebook of Josephine Hopper's, in place of a ledger: one file to a
% notebook, filed under transcripts/notebooks/, the argument the site's slug
% (« battle-of-wash-sq »). Sets the ledger to « notebooks » for everything
% that reads it, so the title block and the watermark behave as they do for a
% batch; the rendering checks the sheets against the notebook's own listing.
\newcommand{\notebook}[1]{\gdef\hopLedger{notebooks}\gdef\hopNotebook{#1}}
\gdef\hopNotebook{}
% The Whitney's accession number for the physical volume — 96.208 … 96.213.
% This is what a citation of the object cites.
\newcommand{\objectnumber}[1]{\gdef\hopObject{#1}}
% The batch this file covers. Leaving it out drops the « (batch N) » from the
% title block rather than printing a number that would be false.
\newcommand{\batch}[1]{\gdef\hopBatch{#1}}
\newcommand{\sheets}[2]{\gdef\hopFirst{#1}\gdef\hopLast{#2}}

% The Whitney's date range for the volume, copied verbatim from its resource
% record — never inferred here, never narrowed to the years this batch
% happens to cover.
\newcommand{\dating}[1]{\gdef\hopDating{#1}}

% The legal notice, printed in the title block and repeated as a diagonal
% watermark on every page of the PDF.
%
% This is not boilerplate and it is not optional. The ledgers are © Heirs of
% Josephine N. Hopper, licensed by the Artists Rights Society; the Whitney
% records the object rights as transferred to the Museum. A transcription
% reproduces Josephine Hopper's text. Every file must therefore say what it is
% on its own face, not only in a README that does not travel with the PDF.
% A file without this line gets no watermark, so omitting it is visible in the
% output rather than silent.
\newcommand{\watermark}[1]{\gdef\hopWatermark{#1}}

\gdef\hopLedger{?}\gdef\hopLedgerTitle{}\gdef\hopObject{}\gdef\hopBatch{}%
\gdef\hopFirst{?}\gdef\hopLast{?}\gdef\hopDating{}\gdef\hopWatermark{}

% --- Keywords ---------------------------------------------------------------
%
% Closing a transcription: the vocabulary under which to search for what these
% leaves record — dealers, exhibiting societies, media, places.
% `npm run manifest` extracts them to tag the whole ledger. This line is the
% single source of the tags; there is no tags file, so no tag can describe
% leaves nobody has read.
%
% It is the one thing in the file that is not on the paper, and it earns its
% place: a reader looking for Keppel needs some way in, and a tag written by
% whoever read the sheets is the only kind that cannot describe unread ones.
\newcommand{\keywords}[1]{\par\smallskip\noindent
  {\footnotesize\textsc{Keywords} --- \textit{#1}}\par}

% --- The anchor -------------------------------------------------------------

% The sheet begins here. Two arguments, and both are needed.
%
% #1 is the Whitney's **resource ref** — 16853 — which is the sheet's whole
% address: it is what the image URL is built from and what turns the facsimile
% as the transcript is scrolled. #2 is the number **written on the paper**, or
% empty where nobody wrote one, which is the case for every cover, flyleaf,
% index and loose insertion — about one sheet in nine.
%
% Keeping both is the whole of the numbering problem in this archive. Three
% sequences overlap: the Hoppers' own leaf numbers, the Whitney's order of
% photographs, and the resource refs, which are in no order at all. Only the
% first is citable and only the third is addressable, so a transcription that
% recorded one of them would be either uncitable or unopenable.
%
% `scripts/render.mjs` checks #2 against what `scripts/catalogue.mjs` parsed
% out of the Whitney's descriptor and fails on a disagreement, because a
% mistyped leaf number silently moves an entry to another year.
\newcommand{\sheet}[2]{%
  \marginnote{\footnotesize\color{gray!70}\textsf{\ifx\relax#2\relax---\else#2\fi}}%
  \label{sheet:#1}\ignorespaces}

% --- The critical apparatus -------------------------------------------------

% Illegible: never guessed. On a ledger leaf this is most often a figure, and
% a guessed figure is a fabricated transaction.
\newcommand{\ill}{\textcolor{red!60!black}{[\,\ldots\,]}}

% A doubtful reading: the word or figure is given, and flagged as uncertain.
%
% Underlining belongs to this macro and to nothing else. Jo Hopper underlines
% her own headings --- « Etchings », « Etching Plates », « Early oil » --- and
% those are set with \emph{} instead, because a reader must be able to tell a
% doubtful reading from her emphasis at a glance and the two cannot both be
% underlines. Which of them keeps the underline is decided by consequence: an
% emphasis shown as italic loses nothing, and a doubtful figure shown as a
% certain one loses everything.
\newcommand{\uncertain}[1]{\uline{#1}}

% An editorial addition — an expanded abbreviation, an implied dollar sign.
\newcommand{\add}[1]{\textcolor{blue!50!black}{[#1]}}

% Struck out in the book. Jo Hopper struck a great deal: a price revised, a
% buyer who withdrew, an exhibition that fell through. What she crossed out is
% part of the record and often the more interesting half of it.
\newcommand{\struck}[1]{\sout{#1}}

% The ink it is written in, where the ink says something.
%
% Book IV is a state machine and the colour is its status field: charges in
% black, receipts in red, the sums in pencil, a rule that holds from leaf 7 to
% the end of the volume. A transcription that records the words and the
% figures and nothing about their colour throws that field away, and the
% accounts then have to guess a receipt from the word « rec'd » --- which the
% volume stops writing on leaf 157. So the ink is recorded, per cell, as what
% is on the paper and not as what it means: \ink{red}{rec'd by check} says
% the words are red, and it is `scripts/accounts.mjs` that says red is money
% received. Black is the default and is not marked. The inks are the three
% the volume uses besides black; a fourth would be added here and in
% `scripts/render.mjs` and `scripts/tei.mjs` in the same commit.
% The file is \input, not \usepackage'd, so @ is not a letter here without saying so.
\makeatletter
\@namedef{hop@ink@red}{red!70!black}
\@namedef{hop@ink@pencil}{gray!55!black}
\@namedef{hop@ink@blue}{blue!55!black}
\newcommand{\ink}[2]{%
  \ifcsname hop@ink@#1\endcsname
    \textcolor{\csname hop@ink@#1\endcsname}{#2}%
  \else
    \PackageError{hopper}{\string\ink{#1}: unknown ink}{The inks are red, pencil and blue.}%
  \fi}
\makeatother

% A rule drawn across the money column above this row's figure.
%
% The writer rules off a run of items and writes their sum under the rule,
% and the rule is the whole of what says « this figure is a sum ». It stands
% at the head of the row the sum is on: \ruledoff & & 4355 & 00. On paper it
% is a line under the last two columns of the four a Book IV leaf has; if
% another volume ever rules a different column, the macro grows an argument.
\newcommand{\ruledoff}{\cline{3-4}}

% A diary entry begins. #1 is the date as she wrote it — « Mar. 12 »,
% « Sunday » — and #2 the date the transcriber assigns it, as a year, a month
% and a day (1947-03-12), or as much of one as the leaf allows. The first is
% what is on the paper; the second is the one editorial claim a diary needs
% that a ledger does not, because a ledger's leaf is its own address and an
% entry's date is its subject. Both are kept, apart, so the claim can be
% argued with. Used only in the notebooks; a ledger has no entries.
\newcommand{\entry}[2]{\par\medskip\noindent
  {\bfseries #1}\ifx\relax#2\relax\else\hspace{0.6em}{\footnotesize\color{gray!60!black}\textsf{#2}}\fi
  \par\nopagebreak\smallskip}

% The transcriber's note. Ours, not theirs.
\newcommand{\note}[1]{\par\smallskip{\small\color{gray!60!black}\textit{[#1]}}\par\smallskip}

% A marginal note **of theirs**, distinct from the body — and distinct from
% \note{} for exactly the reason \note{} exists.
\newcommand{\marginal}[1]{\par\smallskip\noindent\hspace{1em}%
  {\small\color{gray!60!black}#1}\par\smallskip}

% --- What is particular to this archive -------------------------------------

% Whose hand wrote it.
%
% This macro has no counterpart in the Grothendieck preamble, and it is the
% single most important thing in this one. These books were written by two
% people and annotated by more. Edward Hopper drew the work and wrote its
% title and size; Josephine Nivison Hopper wrote everything else — the
% descriptions, the anecdotes about people he refused to discuss, the prices,
% the buyers, the columns; and later hands, curatorial and unidentified, added
% to the leaves after 1967.
%
% A transcription that flattens the three has destroyed the archive's chief
% documentary value, which is that it records *two* working lives and the
% division of labour between them. #1 is `edward`, `jo`, `later` or
% `unidentified`; `unidentified` is the right answer far more often than a
% transcriber's confidence suggests, and it is not a failure to use it.
\newcommand{\hand}[2]{%
  \ifx\relax#1\relax#2\else
    \textcolor{gray!55!black}{\footnotesize\textsf{[#1]}}\,#2%
  \fi}

% Edward Hopper's ink record sketch stands here.
%
% The argument is **what is inscribed on or beside the sketch** — a title, a
% dimension, a note — and never a description of the image. Describing the
% drawing would be writing a new document: the sketch is on the site, one pane
% away, at the resolution the Whitney publishes, and prose about it competes
% with looking at it. Leave the argument empty where nothing is written.
% The mark is drawn rather than set from a symbol font: there is no
% mathematics anywhere in this archive, so no maths package is loaded, and
% $\square$ would be an undefined control sequence — which is exactly how it
% failed the first time the PDFs were compiled.
\newcommand{\sketchmark}{\raisebox{0.15ex}{\color{gray!50!black}%
  \tikz[baseline]{\draw[line width=0.4pt] (0,0) rectangle (1.1ex,1.1ex);}}}
\newcommand{\sketch}[1]{\par\smallskip\noindent
  {\small\color{gray!50!black}\sketchmark~\textsf{ink sketch}%
   \ifx\relax#1\relax\else\ --- #1\fi}\par\smallskip}

% Something printed and pasted to the leaf — a reproduction cut from a
% catalogue, a review, a dealer's label. The argument transcribes its printed
% text. These are the only places in the archive where the words are nobody's
% handwriting, and they date the leaf, which handwriting does not.
\newcommand{\clipping}[1]{\par\smallskip\noindent
  {\small\color{gray!50!black}\textsf{[clipping]}\ #1}\par\smallskip}

% A work's block opens here, under the title **exactly as the leaf gives it** —
% « Les Deux Pigeon » and « Les Poillus » stay misspelled, because the
% misspelling is Hopper's and is how the entry is found.
\newcommand{\work}[1]{\par\medskip\noindent{\bfseries #1}\par\nopagebreak\smallskip}

% --- The ruled table --------------------------------------------------------
%
% Jo Hopper ruled these columns herself and kept them for forty years; the
% table is not a presentation of the content, it *is* the content. So it is
% transcribed as a table and rendered as one, on screen and on paper, from one
% source.
%
% #1 is the column specification; #2 is the header row, `&`-separated like any
% other. A cell she left blank is left blank; a cell that cannot be read takes
% \ill{}.
%
% **Use `Y{...}` for any column that can hold a sentence** — an exhibition, a
% dealer's terms, a note — and `l`, `r` or `c` only for dates, figures and
% short names. See the note on \newcolumntype{Y} above for why that is not a
% matter of taste.
%
% The header row is emitted **bare**, between rules, and is deliberately not
% wrapped in \textbf{}. It cannot be: an `&` inside a braced group is not an
% alignment tab, so `\textbf{Date & Exhibitions}` fails at the first header
% with « Forbidden control sequence found while scanning use of
% \check@nocorr@ » — which is TeX saying, at some distance, that the tabs are
% in the wrong place. Rules carry the header instead, which is how a scholarly
% table sets one anyway. (The reading view has no such constraint and renders
% the same row as <th>.)
\newenvironment{ledgertable}[2]
  {\par\smallskip\small
   \begin{longtable}{#1}
   \hline
   #2 \\
   \hline
   \endhead}
  {\end{longtable}\par\smallskip}

% --- The appendix of notes --------------------------------------------------
%
% Everything above the appendix is on the paper. Nothing in it is.
%
% A note says what a named source says about a work these leaves record --- what
% the holding museum measures the canvas at, which other leaf carries the other
% half of a transaction --- and it is the only prose in the file that a reader
% cannot check against the photograph. On the site that is handled by putting
% the note beside a link to its source; on paper a link is not enough, because
% a PDF outlives the site it was downloaded from and travels away from it. So
% every claim here prints its source in full, name and address both, and every
% note prints the day it was written and what wrote it.
%
% The appendix is assembled by `scripts/pdf.mjs` from
% `src/content/work-notes.json` at compile time and is in no `.tex` file. That
% is the point of it. The transcription is the edition; a note is a later and
% weaker kind of writing *about* the edition, and keeping the two in one source
% file would eventually put a sentence somebody inferred inside the document
% whose whole claim on a reader is that nothing in it was inferred.
%
% Only the notes on works this batch's leaves name are printed, so the appendix
% follows the batch rather than repeating the archive, and a batch that has
% earned no notes yet gets no appendix at all --- which is a true statement
% about how much has been checked, and the same statement `work-notes.json`
% makes by being nearly empty.

% A URL is the one string in this archive that can run a hundred characters
% without a space in it, and an overfull box fails the build here because in a
% transcription an overfull box hides content. So the appendix tells TeX where
% a URL may break and sets its apparatus ragged right, rather than leaving the
% line to overrun in a file nobody recompiles.
\def\UrlBreaks{\do\/\do\-\do\.\do\_\do\?\do\&\do\=\do\+}
\Urlmuskip=0mu plus 1mu

\newcommand{\notesappendix}{\clearpage
  \par\noindent{\large\bfseries Notes on the works}\par\smallskip
  \noindent{\footnotesize\color{gray!60!black}\raggedright
    Not on the paper, and not part of the transcription. Each note below
    concerns a work named on a leaf of this batch and says what a named source
    says about it --- a museum's own catalogue record, or another leaf of these
    ledgers. Every sentence carries the source that states it, printed with its
    address so that it can be followed from a copy of this file alone. Where a
    ledger and a museum disagree the disagreement is printed rather than
    resolved.\par}\par\medskip}

% One work: the title as the leaves give it, then the note.
\newcommand{\worknote}[2]{\par\medskip\noindent{\bfseries #1}\par\nopagebreak
  \smallskip\noindent #2\par}

% One claim under it: what is asserted, and who asserts it.
\newcommand{\workclaim}[2]{\par\smallskip\noindent
  {\footnotesize\color{gray!60!black}\raggedright\hangindent=1.6em\hangafter=0
   \hspace{1em}#1\ --- #2\par}}

% Who wrote the note and when, closing the block. A note is a model's reading
% of retrieved records and is weaker evidence than anything above it; the file
% records that rather than letting the appendix pass for the edition.
\newcommand{\worknoteby}[2]{\par\smallskip\noindent
  {\scriptsize\color{gray!55!black}\hspace{1em}Written #1 by #2.\par}\par\smallskip}

% --- Title ------------------------------------------------------------------

% The watermark is drawn on the shipout background so it sits under the text of
% every page, diagonal and light enough to leave the record readable.
\AddToHook{shipout/background}{%
  \ifx\hopWatermark\empty\else
    \begin{tikzpicture}[remember picture, overlay]
      \node[rotate=52, scale=3.4, align=center, text=gray!18,
            font=\sffamily\bfseries]
        at (current page.center) {\hopWatermark};
    \end{tikzpicture}%
  \fi}

\AtBeginDocument{%
  \begin{center}
    {\large\bfseries Edward and Josephine Hopper --- \hopLedgerTitle}\\[2pt]
    \ifx\hopObject\empty\else
      {\small Whitney Museum of American Art, \hopObject}\\[4pt]
    \fi
    {\small sheets \hopFirst--\hopLast
      \ifx\hopBatch\empty\else\ (batch \hopBatch)\fi}\\[2pt]
    \ifx\hopDating\empty\else
      {\small\color{gray!60!black}The Whitney's dating: \hopDating}\\[2pt]
    \fi
    \ifx\hopWatermark\empty\else
      {\small\bfseries\color{red!45!gray}\hopWatermark}\\[2pt]
    \fi
    {\footnotesize\color{gray!60!black}Transcribed from the digitised sheets published by the
     Whitney Museum of American Art.\\
     \textcopyright{} Heirs of Josephine N. Hopper, licensed by Artists Rights Society (ARS),
     New York.\\
     Uncertain readings are underlined, illegible passages marked [\,\ldots\,],
     editorial additions bracketed.}
  \end{center}
  \vspace{1em}}

\endinput
 rather than \usepackage, and \ProvidesPackage in
% an \input-ed file makes LaTeX warn on every compile that it "requested
% package `'" and got `hopper'. A warning nobody can act on is a warning
% everybody learns to scroll past, which is how the real ones get missed.
\ProvidesFile{hopper.sty}[2026/09/05 Transcriptions of the Hopper artist's ledgers]

\usepackage{array}
% longtable, not tabularx: a leaf of thirty-one exhibition lines must break
% across pages, and tabularx cannot be wrapped in an environment of our own —
% it scans the source for a literal \end{tabularx} and dies on \end{ledgertable}
% with « File ended while scanning use of \TX@get@body ». longtable has no such
% scan, and the wrapping column type below does the job tabularx's X would.
\usepackage{longtable}

% A wrapping column, `Y{0.3}` being three tenths of the measure.
%
% Prose columns must be declared this way and not as `l`. A table of `l`
% columns is set to its natural width, which can be wider than the page, and
% TeX issues **no warning at all** because nothing ever asked the row to fit:
% the first compile of Book I produced a page whose right-hand column ran off
% the paper with a clean log. A fixed-width column wraps, and if it still
% cannot fit, TeX says so — which is what `npm run pdf` then refuses to build
% over.
%
% Leave about 0.03 of the measure per column for the inter-column space: five
% columns can share 0.85, six can share 0.80. Getting it wrong is not a
% guessing game — `npm run pdf` reports the overfull box and how many points
% too wide it is.
\newcolumntype{Y}[1]{>{\raggedright\arraybackslash}p{#1\linewidth}}
\usepackage[english]{babel}
\usepackage{fontspec}
\usepackage{geometry}
\geometry{a4paper,margin=25mm,marginparwidth=28mm}
\usepackage{marginnote}
\usepackage{xcolor}
\usepackage{tikz}
% ulem, not soul. `soul`'s \st and \ul are built on a character-by-character
% scan that XeLaTeX's font machinery breaks: the document fails with an
% undefined control sequence rather than degrading. `normalem` stops it
% hijacking \emph, which the transcriptions use for underlining of Jo Hopper's
% own.
\usepackage[normalem]{ulem}
% ulem sets each underlined word in a box TeX can neither hyphenate nor break
% after, so a paragraph dense in \uncertain{} readings can reach a state where
% no legal breakpoint exists. \emergencystretch lets TeX loosen it on a third
% pass instead of overrunning: an overfull line in a transcription hides
% content.
\setlength{\emergencystretch}{3em}
\usepackage{hyperref}
\hypersetup{colorlinks=true,linkcolor=black,urlcolor=black,pdfborder={0 0 0}}

% --- Batch metadata ---------------------------------------------------------
% Taken as they stand by the rendering script, which shows them at the head of
% the reading view. They fix what the file claims to cover: without them a
% transcription floats unmoored in an archive of five hundred sheets.

\newcommand{\ledger}[1]{\gdef\hopLedger{#1}}
\newcommand{\ledgertitle}[1]{\gdef\hopLedgerTitle{#1}}
% A notebook of Josephine Hopper's, in place of a ledger: one file to a
% notebook, filed under transcripts/notebooks/, the argument the site's slug
% (« battle-of-wash-sq »). Sets the ledger to « notebooks » for everything
% that reads it, so the title block and the watermark behave as they do for a
% batch; the rendering checks the sheets against the notebook's own listing.
\newcommand{\notebook}[1]{\gdef\hopLedger{notebooks}\gdef\hopNotebook{#1}}
\gdef\hopNotebook{}
% The Whitney's accession number for the physical volume — 96.208 … 96.213.
% This is what a citation of the object cites.
\newcommand{\objectnumber}[1]{\gdef\hopObject{#1}}
% The batch this file covers. Leaving it out drops the « (batch N) » from the
% title block rather than printing a number that would be false.
\newcommand{\batch}[1]{\gdef\hopBatch{#1}}
\newcommand{\sheets}[2]{\gdef\hopFirst{#1}\gdef\hopLast{#2}}

% The Whitney's date range for the volume, copied verbatim from its resource
% record — never inferred here, never narrowed to the years this batch
% happens to cover.
\newcommand{\dating}[1]{\gdef\hopDating{#1}}

% The legal notice, printed in the title block and repeated as a diagonal
% watermark on every page of the PDF.
%
% This is not boilerplate and it is not optional. The ledgers are © Heirs of
% Josephine N. Hopper, licensed by the Artists Rights Society; the Whitney
% records the object rights as transferred to the Museum. A transcription
% reproduces Josephine Hopper's text. Every file must therefore say what it is
% on its own face, not only in a README that does not travel with the PDF.
% A file without this line gets no watermark, so omitting it is visible in the
% output rather than silent.
\newcommand{\watermark}[1]{\gdef\hopWatermark{#1}}

\gdef\hopLedger{?}\gdef\hopLedgerTitle{}\gdef\hopObject{}\gdef\hopBatch{}%
\gdef\hopFirst{?}\gdef\hopLast{?}\gdef\hopDating{}\gdef\hopWatermark{}

% --- Keywords ---------------------------------------------------------------
%
% Closing a transcription: the vocabulary under which to search for what these
% leaves record — dealers, exhibiting societies, media, places.
% `npm run manifest` extracts them to tag the whole ledger. This line is the
% single source of the tags; there is no tags file, so no tag can describe
% leaves nobody has read.
%
% It is the one thing in the file that is not on the paper, and it earns its
% place: a reader looking for Keppel needs some way in, and a tag written by
% whoever read the sheets is the only kind that cannot describe unread ones.
\newcommand{\keywords}[1]{\par\smallskip\noindent
  {\footnotesize\textsc{Keywords} --- \textit{#1}}\par}

% --- The anchor -------------------------------------------------------------

% The sheet begins here. Two arguments, and both are needed.
%
% #1 is the Whitney's **resource ref** — 16853 — which is the sheet's whole
% address: it is what the image URL is built from and what turns the facsimile
% as the transcript is scrolled. #2 is the number **written on the paper**, or
% empty where nobody wrote one, which is the case for every cover, flyleaf,
% index and loose insertion — about one sheet in nine.
%
% Keeping both is the whole of the numbering problem in this archive. Three
% sequences overlap: the Hoppers' own leaf numbers, the Whitney's order of
% photographs, and the resource refs, which are in no order at all. Only the
% first is citable and only the third is addressable, so a transcription that
% recorded one of them would be either uncitable or unopenable.
%
% `scripts/render.mjs` checks #2 against what `scripts/catalogue.mjs` parsed
% out of the Whitney's descriptor and fails on a disagreement, because a
% mistyped leaf number silently moves an entry to another year.
\newcommand{\sheet}[2]{%
  \marginnote{\footnotesize\color{gray!70}\textsf{\ifx\relax#2\relax---\else#2\fi}}%
  \label{sheet:#1}\ignorespaces}

% --- The critical apparatus -------------------------------------------------

% Illegible: never guessed. On a ledger leaf this is most often a figure, and
% a guessed figure is a fabricated transaction.
\newcommand{\ill}{\textcolor{red!60!black}{[\,\ldots\,]}}

% A doubtful reading: the word or figure is given, and flagged as uncertain.
%
% Underlining belongs to this macro and to nothing else. Jo Hopper underlines
% her own headings --- « Etchings », « Etching Plates », « Early oil » --- and
% those are set with \emph{} instead, because a reader must be able to tell a
% doubtful reading from her emphasis at a glance and the two cannot both be
% underlines. Which of them keeps the underline is decided by consequence: an
% emphasis shown as italic loses nothing, and a doubtful figure shown as a
% certain one loses everything.
\newcommand{\uncertain}[1]{\uline{#1}}

% An editorial addition — an expanded abbreviation, an implied dollar sign.
\newcommand{\add}[1]{\textcolor{blue!50!black}{[#1]}}

% Struck out in the book. Jo Hopper struck a great deal: a price revised, a
% buyer who withdrew, an exhibition that fell through. What she crossed out is
% part of the record and often the more interesting half of it.
\newcommand{\struck}[1]{\sout{#1}}

% The ink it is written in, where the ink says something.
%
% Book IV is a state machine and the colour is its status field: charges in
% black, receipts in red, the sums in pencil, a rule that holds from leaf 7 to
% the end of the volume. A transcription that records the words and the
% figures and nothing about their colour throws that field away, and the
% accounts then have to guess a receipt from the word « rec'd » --- which the
% volume stops writing on leaf 157. So the ink is recorded, per cell, as what
% is on the paper and not as what it means: \ink{red}{rec'd by check} says
% the words are red, and it is `scripts/accounts.mjs` that says red is money
% received. Black is the default and is not marked. The inks are the three
% the volume uses besides black; a fourth would be added here and in
% `scripts/render.mjs` and `scripts/tei.mjs` in the same commit.
% The file is \input, not \usepackage'd, so @ is not a letter here without saying so.
\makeatletter
\@namedef{hop@ink@red}{red!70!black}
\@namedef{hop@ink@pencil}{gray!55!black}
\@namedef{hop@ink@blue}{blue!55!black}
\newcommand{\ink}[2]{%
  \ifcsname hop@ink@#1\endcsname
    \textcolor{\csname hop@ink@#1\endcsname}{#2}%
  \else
    \PackageError{hopper}{\string\ink{#1}: unknown ink}{The inks are red, pencil and blue.}%
  \fi}
\makeatother

% A rule drawn across the money column above this row's figure.
%
% The writer rules off a run of items and writes their sum under the rule,
% and the rule is the whole of what says « this figure is a sum ». It stands
% at the head of the row the sum is on: \ruledoff & & 4355 & 00. On paper it
% is a line under the last two columns of the four a Book IV leaf has; if
% another volume ever rules a different column, the macro grows an argument.
\newcommand{\ruledoff}{\cline{3-4}}

% A diary entry begins. #1 is the date as she wrote it — « Mar. 12 »,
% « Sunday » — and #2 the date the transcriber assigns it, as a year, a month
% and a day (1947-03-12), or as much of one as the leaf allows. The first is
% what is on the paper; the second is the one editorial claim a diary needs
% that a ledger does not, because a ledger's leaf is its own address and an
% entry's date is its subject. Both are kept, apart, so the claim can be
% argued with. Used only in the notebooks; a ledger has no entries.
\newcommand{\entry}[2]{\par\medskip\noindent
  {\bfseries #1}\ifx\relax#2\relax\else\hspace{0.6em}{\footnotesize\color{gray!60!black}\textsf{#2}}\fi
  \par\nopagebreak\smallskip}

% The transcriber's note. Ours, not theirs.
\newcommand{\note}[1]{\par\smallskip{\small\color{gray!60!black}\textit{[#1]}}\par\smallskip}

% A marginal note **of theirs**, distinct from the body — and distinct from
% \note{} for exactly the reason \note{} exists.
\newcommand{\marginal}[1]{\par\smallskip\noindent\hspace{1em}%
  {\small\color{gray!60!black}#1}\par\smallskip}

% --- What is particular to this archive -------------------------------------

% Whose hand wrote it.
%
% This macro has no counterpart in the Grothendieck preamble, and it is the
% single most important thing in this one. These books were written by two
% people and annotated by more. Edward Hopper drew the work and wrote its
% title and size; Josephine Nivison Hopper wrote everything else — the
% descriptions, the anecdotes about people he refused to discuss, the prices,
% the buyers, the columns; and later hands, curatorial and unidentified, added
% to the leaves after 1967.
%
% A transcription that flattens the three has destroyed the archive's chief
% documentary value, which is that it records *two* working lives and the
% division of labour between them. #1 is `edward`, `jo`, `later` or
% `unidentified`; `unidentified` is the right answer far more often than a
% transcriber's confidence suggests, and it is not a failure to use it.
\newcommand{\hand}[2]{%
  \ifx\relax#1\relax#2\else
    \textcolor{gray!55!black}{\footnotesize\textsf{[#1]}}\,#2%
  \fi}

% Edward Hopper's ink record sketch stands here.
%
% The argument is **what is inscribed on or beside the sketch** — a title, a
% dimension, a note — and never a description of the image. Describing the
% drawing would be writing a new document: the sketch is on the site, one pane
% away, at the resolution the Whitney publishes, and prose about it competes
% with looking at it. Leave the argument empty where nothing is written.
% The mark is drawn rather than set from a symbol font: there is no
% mathematics anywhere in this archive, so no maths package is loaded, and
% $\square$ would be an undefined control sequence — which is exactly how it
% failed the first time the PDFs were compiled.
\newcommand{\sketchmark}{\raisebox{0.15ex}{\color{gray!50!black}%
  \tikz[baseline]{\draw[line width=0.4pt] (0,0) rectangle (1.1ex,1.1ex);}}}
\newcommand{\sketch}[1]{\par\smallskip\noindent
  {\small\color{gray!50!black}\sketchmark~\textsf{ink sketch}%
   \ifx\relax#1\relax\else\ --- #1\fi}\par\smallskip}

% Something printed and pasted to the leaf — a reproduction cut from a
% catalogue, a review, a dealer's label. The argument transcribes its printed
% text. These are the only places in the archive where the words are nobody's
% handwriting, and they date the leaf, which handwriting does not.
\newcommand{\clipping}[1]{\par\smallskip\noindent
  {\small\color{gray!50!black}\textsf{[clipping]}\ #1}\par\smallskip}

% A work's block opens here, under the title **exactly as the leaf gives it** —
% « Les Deux Pigeon » and « Les Poillus » stay misspelled, because the
% misspelling is Hopper's and is how the entry is found.
\newcommand{\work}[1]{\par\medskip\noindent{\bfseries #1}\par\nopagebreak\smallskip}

% --- The ruled table --------------------------------------------------------
%
% Jo Hopper ruled these columns herself and kept them for forty years; the
% table is not a presentation of the content, it *is* the content. So it is
% transcribed as a table and rendered as one, on screen and on paper, from one
% source.
%
% #1 is the column specification; #2 is the header row, `&`-separated like any
% other. A cell she left blank is left blank; a cell that cannot be read takes
% \ill{}.
%
% **Use `Y{...}` for any column that can hold a sentence** — an exhibition, a
% dealer's terms, a note — and `l`, `r` or `c` only for dates, figures and
% short names. See the note on \newcolumntype{Y} above for why that is not a
% matter of taste.
%
% The header row is emitted **bare**, between rules, and is deliberately not
% wrapped in \textbf{}. It cannot be: an `&` inside a braced group is not an
% alignment tab, so `\textbf{Date & Exhibitions}` fails at the first header
% with « Forbidden control sequence found while scanning use of
% \check@nocorr@ » — which is TeX saying, at some distance, that the tabs are
% in the wrong place. Rules carry the header instead, which is how a scholarly
% table sets one anyway. (The reading view has no such constraint and renders
% the same row as <th>.)
\newenvironment{ledgertable}[2]
  {\par\smallskip\small
   \begin{longtable}{#1}
   \hline
   #2 \\
   \hline
   \endhead}
  {\end{longtable}\par\smallskip}

% --- The appendix of notes --------------------------------------------------
%
% Everything above the appendix is on the paper. Nothing in it is.
%
% A note says what a named source says about a work these leaves record --- what
% the holding museum measures the canvas at, which other leaf carries the other
% half of a transaction --- and it is the only prose in the file that a reader
% cannot check against the photograph. On the site that is handled by putting
% the note beside a link to its source; on paper a link is not enough, because
% a PDF outlives the site it was downloaded from and travels away from it. So
% every claim here prints its source in full, name and address both, and every
% note prints the day it was written and what wrote it.
%
% The appendix is assembled by `scripts/pdf.mjs` from
% `src/content/work-notes.json` at compile time and is in no `.tex` file. That
% is the point of it. The transcription is the edition; a note is a later and
% weaker kind of writing *about* the edition, and keeping the two in one source
% file would eventually put a sentence somebody inferred inside the document
% whose whole claim on a reader is that nothing in it was inferred.
%
% Only the notes on works this batch's leaves name are printed, so the appendix
% follows the batch rather than repeating the archive, and a batch that has
% earned no notes yet gets no appendix at all --- which is a true statement
% about how much has been checked, and the same statement `work-notes.json`
% makes by being nearly empty.

% A URL is the one string in this archive that can run a hundred characters
% without a space in it, and an overfull box fails the build here because in a
% transcription an overfull box hides content. So the appendix tells TeX where
% a URL may break and sets its apparatus ragged right, rather than leaving the
% line to overrun in a file nobody recompiles.
\def\UrlBreaks{\do\/\do\-\do\.\do\_\do\?\do\&\do\=\do\+}
\Urlmuskip=0mu plus 1mu

\newcommand{\notesappendix}{\clearpage
  \par\noindent{\large\bfseries Notes on the works}\par\smallskip
  \noindent{\footnotesize\color{gray!60!black}\raggedright
    Not on the paper, and not part of the transcription. Each note below
    concerns a work named on a leaf of this batch and says what a named source
    says about it --- a museum's own catalogue record, or another leaf of these
    ledgers. Every sentence carries the source that states it, printed with its
    address so that it can be followed from a copy of this file alone. Where a
    ledger and a museum disagree the disagreement is printed rather than
    resolved.\par}\par\medskip}

% One work: the title as the leaves give it, then the note.
\newcommand{\worknote}[2]{\par\medskip\noindent{\bfseries #1}\par\nopagebreak
  \smallskip\noindent #2\par}

% One claim under it: what is asserted, and who asserts it.
\newcommand{\workclaim}[2]{\par\smallskip\noindent
  {\footnotesize\color{gray!60!black}\raggedright\hangindent=1.6em\hangafter=0
   \hspace{1em}#1\ --- #2\par}}

% Who wrote the note and when, closing the block. A note is a model's reading
% of retrieved records and is weaker evidence than anything above it; the file
% records that rather than letting the appendix pass for the edition.
\newcommand{\worknoteby}[2]{\par\smallskip\noindent
  {\scriptsize\color{gray!55!black}\hspace{1em}Written #1 by #2.\par}\par\smallskip}

% --- Title ------------------------------------------------------------------

% The watermark is drawn on the shipout background so it sits under the text of
% every page, diagonal and light enough to leave the record readable.
\AddToHook{shipout/background}{%
  \ifx\hopWatermark\empty\else
    \begin{tikzpicture}[remember picture, overlay]
      \node[rotate=52, scale=3.4, align=center, text=gray!18,
            font=\sffamily\bfseries]
        at (current page.center) {\hopWatermark};
    \end{tikzpicture}%
  \fi}

\AtBeginDocument{%
  \begin{center}
    {\large\bfseries Edward and Josephine Hopper --- \hopLedgerTitle}\\[2pt]
    \ifx\hopObject\empty\else
      {\small Whitney Museum of American Art, \hopObject}\\[4pt]
    \fi
    {\small sheets \hopFirst--\hopLast
      \ifx\hopBatch\empty\else\ (batch \hopBatch)\fi}\\[2pt]
    \ifx\hopDating\empty\else
      {\small\color{gray!60!black}The Whitney's dating: \hopDating}\\[2pt]
    \fi
    \ifx\hopWatermark\empty\else
      {\small\bfseries\color{red!45!gray}\hopWatermark}\\[2pt]
    \fi
    {\footnotesize\color{gray!60!black}Transcribed from the digitised sheets published by the
     Whitney Museum of American Art.\\
     \textcopyright{} Heirs of Josephine N. Hopper, licensed by Artists Rights Society (ARS),
     New York.\\
     Uncertain readings are underlined, illegible passages marked [\,\ldots\,],
     editorial additions bracketed.}
  \end{center}
  \vspace{1em}}

\endinput
 rather than \usepackage, and \ProvidesPackage in
% an \input-ed file makes LaTeX warn on every compile that it "requested
% package `'" and got `hopper'. A warning nobody can act on is a warning
% everybody learns to scroll past, which is how the real ones get missed.
\ProvidesFile{hopper.sty}[2026/09/05 Transcriptions of the Hopper artist's ledgers]

\usepackage{array}
% longtable, not tabularx: a leaf of thirty-one exhibition lines must break
% across pages, and tabularx cannot be wrapped in an environment of our own —
% it scans the source for a literal \end{tabularx} and dies on \end{ledgertable}
% with « File ended while scanning use of \TX@get@body ». longtable has no such
% scan, and the wrapping column type below does the job tabularx's X would.
\usepackage{longtable}

% A wrapping column, `Y{0.3}` being three tenths of the measure.
%
% Prose columns must be declared this way and not as `l`. A table of `l`
% columns is set to its natural width, which can be wider than the page, and
% TeX issues **no warning at all** because nothing ever asked the row to fit:
% the first compile of Book I produced a page whose right-hand column ran off
% the paper with a clean log. A fixed-width column wraps, and if it still
% cannot fit, TeX says so — which is what `npm run pdf` then refuses to build
% over.
%
% Leave about 0.03 of the measure per column for the inter-column space: five
% columns can share 0.85, six can share 0.80. Getting it wrong is not a
% guessing game — `npm run pdf` reports the overfull box and how many points
% too wide it is.
\newcolumntype{Y}[1]{>{\raggedright\arraybackslash}p{#1\linewidth}}
\usepackage[english]{babel}
\usepackage{fontspec}
\usepackage{geometry}
\geometry{a4paper,margin=25mm,marginparwidth=28mm}
\usepackage{marginnote}
\usepackage{xcolor}
\usepackage{tikz}
% ulem, not soul. `soul`'s \st and \ul are built on a character-by-character
% scan that XeLaTeX's font machinery breaks: the document fails with an
% undefined control sequence rather than degrading. `normalem` stops it
% hijacking \emph, which the transcriptions use for underlining of Jo Hopper's
% own.
\usepackage[normalem]{ulem}
% ulem sets each underlined word in a box TeX can neither hyphenate nor break
% after, so a paragraph dense in \uncertain{} readings can reach a state where
% no legal breakpoint exists. \emergencystretch lets TeX loosen it on a third
% pass instead of overrunning: an overfull line in a transcription hides
% content.
\setlength{\emergencystretch}{3em}
\usepackage{hyperref}
\hypersetup{colorlinks=true,linkcolor=black,urlcolor=black,pdfborder={0 0 0}}

% --- Batch metadata ---------------------------------------------------------
% Taken as they stand by the rendering script, which shows them at the head of
% the reading view. They fix what the file claims to cover: without them a
% transcription floats unmoored in an archive of five hundred sheets.

\newcommand{\ledger}[1]{\gdef\hopLedger{#1}}
\newcommand{\ledgertitle}[1]{\gdef\hopLedgerTitle{#1}}
% A notebook of Josephine Hopper's, in place of a ledger: one file to a
% notebook, filed under transcripts/notebooks/, the argument the site's slug
% (« battle-of-wash-sq »). Sets the ledger to « notebooks » for everything
% that reads it, so the title block and the watermark behave as they do for a
% batch; the rendering checks the sheets against the notebook's own listing.
\newcommand{\notebook}[1]{\gdef\hopLedger{notebooks}\gdef\hopNotebook{#1}}
\gdef\hopNotebook{}
% The Whitney's accession number for the physical volume — 96.208 … 96.213.
% This is what a citation of the object cites.
\newcommand{\objectnumber}[1]{\gdef\hopObject{#1}}
% The batch this file covers. Leaving it out drops the « (batch N) » from the
% title block rather than printing a number that would be false.
\newcommand{\batch}[1]{\gdef\hopBatch{#1}}
\newcommand{\sheets}[2]{\gdef\hopFirst{#1}\gdef\hopLast{#2}}

% The Whitney's date range for the volume, copied verbatim from its resource
% record — never inferred here, never narrowed to the years this batch
% happens to cover.
\newcommand{\dating}[1]{\gdef\hopDating{#1}}

% The legal notice, printed in the title block and repeated as a diagonal
% watermark on every page of the PDF.
%
% This is not boilerplate and it is not optional. The ledgers are © Heirs of
% Josephine N. Hopper, licensed by the Artists Rights Society; the Whitney
% records the object rights as transferred to the Museum. A transcription
% reproduces Josephine Hopper's text. Every file must therefore say what it is
% on its own face, not only in a README that does not travel with the PDF.
% A file without this line gets no watermark, so omitting it is visible in the
% output rather than silent.
\newcommand{\watermark}[1]{\gdef\hopWatermark{#1}}

\gdef\hopLedger{?}\gdef\hopLedgerTitle{}\gdef\hopObject{}\gdef\hopBatch{}%
\gdef\hopFirst{?}\gdef\hopLast{?}\gdef\hopDating{}\gdef\hopWatermark{}

% --- Keywords ---------------------------------------------------------------
%
% Closing a transcription: the vocabulary under which to search for what these
% leaves record — dealers, exhibiting societies, media, places.
% `npm run manifest` extracts them to tag the whole ledger. This line is the
% single source of the tags; there is no tags file, so no tag can describe
% leaves nobody has read.
%
% It is the one thing in the file that is not on the paper, and it earns its
% place: a reader looking for Keppel needs some way in, and a tag written by
% whoever read the sheets is the only kind that cannot describe unread ones.
\newcommand{\keywords}[1]{\par\smallskip\noindent
  {\footnotesize\textsc{Keywords} --- \textit{#1}}\par}

% --- The anchor -------------------------------------------------------------

% The sheet begins here. Two arguments, and both are needed.
%
% #1 is the Whitney's **resource ref** — 16853 — which is the sheet's whole
% address: it is what the image URL is built from and what turns the facsimile
% as the transcript is scrolled. #2 is the number **written on the paper**, or
% empty where nobody wrote one, which is the case for every cover, flyleaf,
% index and loose insertion — about one sheet in nine.
%
% Keeping both is the whole of the numbering problem in this archive. Three
% sequences overlap: the Hoppers' own leaf numbers, the Whitney's order of
% photographs, and the resource refs, which are in no order at all. Only the
% first is citable and only the third is addressable, so a transcription that
% recorded one of them would be either uncitable or unopenable.
%
% `scripts/render.mjs` checks #2 against what `scripts/catalogue.mjs` parsed
% out of the Whitney's descriptor and fails on a disagreement, because a
% mistyped leaf number silently moves an entry to another year.
\newcommand{\sheet}[2]{%
  \marginnote{\footnotesize\color{gray!70}\textsf{\ifx\relax#2\relax---\else#2\fi}}%
  \label{sheet:#1}\ignorespaces}

% --- The critical apparatus -------------------------------------------------

% Illegible: never guessed. On a ledger leaf this is most often a figure, and
% a guessed figure is a fabricated transaction.
\newcommand{\ill}{\textcolor{red!60!black}{[\,\ldots\,]}}

% A doubtful reading: the word or figure is given, and flagged as uncertain.
%
% Underlining belongs to this macro and to nothing else. Jo Hopper underlines
% her own headings --- « Etchings », « Etching Plates », « Early oil » --- and
% those are set with \emph{} instead, because a reader must be able to tell a
% doubtful reading from her emphasis at a glance and the two cannot both be
% underlines. Which of them keeps the underline is decided by consequence: an
% emphasis shown as italic loses nothing, and a doubtful figure shown as a
% certain one loses everything.
\newcommand{\uncertain}[1]{\uline{#1}}

% An editorial addition — an expanded abbreviation, an implied dollar sign.
\newcommand{\add}[1]{\textcolor{blue!50!black}{[#1]}}

% Struck out in the book. Jo Hopper struck a great deal: a price revised, a
% buyer who withdrew, an exhibition that fell through. What she crossed out is
% part of the record and often the more interesting half of it.
\newcommand{\struck}[1]{\sout{#1}}

% The ink it is written in, where the ink says something.
%
% Book IV is a state machine and the colour is its status field: charges in
% black, receipts in red, the sums in pencil, a rule that holds from leaf 7 to
% the end of the volume. A transcription that records the words and the
% figures and nothing about their colour throws that field away, and the
% accounts then have to guess a receipt from the word « rec'd » --- which the
% volume stops writing on leaf 157. So the ink is recorded, per cell, as what
% is on the paper and not as what it means: \ink{red}{rec'd by check} says
% the words are red, and it is `scripts/accounts.mjs` that says red is money
% received. Black is the default and is not marked. The inks are the three
% the volume uses besides black; a fourth would be added here and in
% `scripts/render.mjs` and `scripts/tei.mjs` in the same commit.
% The file is \input, not \usepackage'd, so @ is not a letter here without saying so.
\makeatletter
\@namedef{hop@ink@red}{red!70!black}
\@namedef{hop@ink@pencil}{gray!55!black}
\@namedef{hop@ink@blue}{blue!55!black}
\newcommand{\ink}[2]{%
  \ifcsname hop@ink@#1\endcsname
    \textcolor{\csname hop@ink@#1\endcsname}{#2}%
  \else
    \PackageError{hopper}{\string\ink{#1}: unknown ink}{The inks are red, pencil and blue.}%
  \fi}
\makeatother

% A rule drawn across the money column above this row's figure.
%
% The writer rules off a run of items and writes their sum under the rule,
% and the rule is the whole of what says « this figure is a sum ». It stands
% at the head of the row the sum is on: \ruledoff & & 4355 & 00. On paper it
% is a line under the last two columns of the four a Book IV leaf has; if
% another volume ever rules a different column, the macro grows an argument.
\newcommand{\ruledoff}{\cline{3-4}}

% A diary entry begins. #1 is the date as she wrote it — « Mar. 12 »,
% « Sunday » — and #2 the date the transcriber assigns it, as a year, a month
% and a day (1947-03-12), or as much of one as the leaf allows. The first is
% what is on the paper; the second is the one editorial claim a diary needs
% that a ledger does not, because a ledger's leaf is its own address and an
% entry's date is its subject. Both are kept, apart, so the claim can be
% argued with. Used only in the notebooks; a ledger has no entries.
\newcommand{\entry}[2]{\par\medskip\noindent
  {\bfseries #1}\ifx\relax#2\relax\else\hspace{0.6em}{\footnotesize\color{gray!60!black}\textsf{#2}}\fi
  \par\nopagebreak\smallskip}

% The transcriber's note. Ours, not theirs.
\newcommand{\note}[1]{\par\smallskip{\small\color{gray!60!black}\textit{[#1]}}\par\smallskip}

% A marginal note **of theirs**, distinct from the body — and distinct from
% \note{} for exactly the reason \note{} exists.
\newcommand{\marginal}[1]{\par\smallskip\noindent\hspace{1em}%
  {\small\color{gray!60!black}#1}\par\smallskip}

% --- What is particular to this archive -------------------------------------

% Whose hand wrote it.
%
% This macro has no counterpart in the Grothendieck preamble, and it is the
% single most important thing in this one. These books were written by two
% people and annotated by more. Edward Hopper drew the work and wrote its
% title and size; Josephine Nivison Hopper wrote everything else — the
% descriptions, the anecdotes about people he refused to discuss, the prices,
% the buyers, the columns; and later hands, curatorial and unidentified, added
% to the leaves after 1967.
%
% A transcription that flattens the three has destroyed the archive's chief
% documentary value, which is that it records *two* working lives and the
% division of labour between them. #1 is `edward`, `jo`, `later` or
% `unidentified`; `unidentified` is the right answer far more often than a
% transcriber's confidence suggests, and it is not a failure to use it.
\newcommand{\hand}[2]{%
  \ifx\relax#1\relax#2\else
    \textcolor{gray!55!black}{\footnotesize\textsf{[#1]}}\,#2%
  \fi}

% Edward Hopper's ink record sketch stands here.
%
% The argument is **what is inscribed on or beside the sketch** — a title, a
% dimension, a note — and never a description of the image. Describing the
% drawing would be writing a new document: the sketch is on the site, one pane
% away, at the resolution the Whitney publishes, and prose about it competes
% with looking at it. Leave the argument empty where nothing is written.
% The mark is drawn rather than set from a symbol font: there is no
% mathematics anywhere in this archive, so no maths package is loaded, and
% $\square$ would be an undefined control sequence — which is exactly how it
% failed the first time the PDFs were compiled.
\newcommand{\sketchmark}{\raisebox{0.15ex}{\color{gray!50!black}%
  \tikz[baseline]{\draw[line width=0.4pt] (0,0) rectangle (1.1ex,1.1ex);}}}
\newcommand{\sketch}[1]{\par\smallskip\noindent
  {\small\color{gray!50!black}\sketchmark~\textsf{ink sketch}%
   \ifx\relax#1\relax\else\ --- #1\fi}\par\smallskip}

% Something printed and pasted to the leaf — a reproduction cut from a
% catalogue, a review, a dealer's label. The argument transcribes its printed
% text. These are the only places in the archive where the words are nobody's
% handwriting, and they date the leaf, which handwriting does not.
\newcommand{\clipping}[1]{\par\smallskip\noindent
  {\small\color{gray!50!black}\textsf{[clipping]}\ #1}\par\smallskip}

% A work's block opens here, under the title **exactly as the leaf gives it** —
% « Les Deux Pigeon » and « Les Poillus » stay misspelled, because the
% misspelling is Hopper's and is how the entry is found.
\newcommand{\work}[1]{\par\medskip\noindent{\bfseries #1}\par\nopagebreak\smallskip}

% --- The ruled table --------------------------------------------------------
%
% Jo Hopper ruled these columns herself and kept them for forty years; the
% table is not a presentation of the content, it *is* the content. So it is
% transcribed as a table and rendered as one, on screen and on paper, from one
% source.
%
% #1 is the column specification; #2 is the header row, `&`-separated like any
% other. A cell she left blank is left blank; a cell that cannot be read takes
% \ill{}.
%
% **Use `Y{...}` for any column that can hold a sentence** — an exhibition, a
% dealer's terms, a note — and `l`, `r` or `c` only for dates, figures and
% short names. See the note on \newcolumntype{Y} above for why that is not a
% matter of taste.
%
% The header row is emitted **bare**, between rules, and is deliberately not
% wrapped in \textbf{}. It cannot be: an `&` inside a braced group is not an
% alignment tab, so `\textbf{Date & Exhibitions}` fails at the first header
% with « Forbidden control sequence found while scanning use of
% \check@nocorr@ » — which is TeX saying, at some distance, that the tabs are
% in the wrong place. Rules carry the header instead, which is how a scholarly
% table sets one anyway. (The reading view has no such constraint and renders
% the same row as <th>.)
\newenvironment{ledgertable}[2]
  {\par\smallskip\small
   \begin{longtable}{#1}
   \hline
   #2 \\
   \hline
   \endhead}
  {\end{longtable}\par\smallskip}

% --- The appendix of notes --------------------------------------------------
%
% Everything above the appendix is on the paper. Nothing in it is.
%
% A note says what a named source says about a work these leaves record --- what
% the holding museum measures the canvas at, which other leaf carries the other
% half of a transaction --- and it is the only prose in the file that a reader
% cannot check against the photograph. On the site that is handled by putting
% the note beside a link to its source; on paper a link is not enough, because
% a PDF outlives the site it was downloaded from and travels away from it. So
% every claim here prints its source in full, name and address both, and every
% note prints the day it was written and what wrote it.
%
% The appendix is assembled by `scripts/pdf.mjs` from
% `src/content/work-notes.json` at compile time and is in no `.tex` file. That
% is the point of it. The transcription is the edition; a note is a later and
% weaker kind of writing *about* the edition, and keeping the two in one source
% file would eventually put a sentence somebody inferred inside the document
% whose whole claim on a reader is that nothing in it was inferred.
%
% Only the notes on works this batch's leaves name are printed, so the appendix
% follows the batch rather than repeating the archive, and a batch that has
% earned no notes yet gets no appendix at all --- which is a true statement
% about how much has been checked, and the same statement `work-notes.json`
% makes by being nearly empty.

% A URL is the one string in this archive that can run a hundred characters
% without a space in it, and an overfull box fails the build here because in a
% transcription an overfull box hides content. So the appendix tells TeX where
% a URL may break and sets its apparatus ragged right, rather than leaving the
% line to overrun in a file nobody recompiles.
\def\UrlBreaks{\do\/\do\-\do\.\do\_\do\?\do\&\do\=\do\+}
\Urlmuskip=0mu plus 1mu

\newcommand{\notesappendix}{\clearpage
  \par\noindent{\large\bfseries Notes on the works}\par\smallskip
  \noindent{\footnotesize\color{gray!60!black}\raggedright
    Not on the paper, and not part of the transcription. Each note below
    concerns a work named on a leaf of this batch and says what a named source
    says about it --- a museum's own catalogue record, or another leaf of these
    ledgers. Every sentence carries the source that states it, printed with its
    address so that it can be followed from a copy of this file alone. Where a
    ledger and a museum disagree the disagreement is printed rather than
    resolved.\par}\par\medskip}

% One work: the title as the leaves give it, then the note.
\newcommand{\worknote}[2]{\par\medskip\noindent{\bfseries #1}\par\nopagebreak
  \smallskip\noindent #2\par}

% One claim under it: what is asserted, and who asserts it.
\newcommand{\workclaim}[2]{\par\smallskip\noindent
  {\footnotesize\color{gray!60!black}\raggedright\hangindent=1.6em\hangafter=0
   \hspace{1em}#1\ --- #2\par}}

% Who wrote the note and when, closing the block. A note is a model's reading
% of retrieved records and is weaker evidence than anything above it; the file
% records that rather than letting the appendix pass for the edition.
\newcommand{\worknoteby}[2]{\par\smallskip\noindent
  {\scriptsize\color{gray!55!black}\hspace{1em}Written #1 by #2.\par}\par\smallskip}

% --- Title ------------------------------------------------------------------

% The watermark is drawn on the shipout background so it sits under the text of
% every page, diagonal and light enough to leave the record readable.
\AddToHook{shipout/background}{%
  \ifx\hopWatermark\empty\else
    \begin{tikzpicture}[remember picture, overlay]
      \node[rotate=52, scale=3.4, align=center, text=gray!18,
            font=\sffamily\bfseries]
        at (current page.center) {\hopWatermark};
    \end{tikzpicture}%
  \fi}

\AtBeginDocument{%
  \begin{center}
    {\large\bfseries Edward and Josephine Hopper --- \hopLedgerTitle}\\[2pt]
    \ifx\hopObject\empty\else
      {\small Whitney Museum of American Art, \hopObject}\\[4pt]
    \fi
    {\small sheets \hopFirst--\hopLast
      \ifx\hopBatch\empty\else\ (batch \hopBatch)\fi}\\[2pt]
    \ifx\hopDating\empty\else
      {\small\color{gray!60!black}The Whitney's dating: \hopDating}\\[2pt]
    \fi
    \ifx\hopWatermark\empty\else
      {\small\bfseries\color{red!45!gray}\hopWatermark}\\[2pt]
    \fi
    {\footnotesize\color{gray!60!black}Transcribed from the digitised sheets published by the
     Whitney Museum of American Art.\\
     \textcopyright{} Heirs of Josephine N. Hopper, licensed by Artists Rights Society (ARS),
     New York.\\
     Uncertain readings are underlined, illegible passages marked [\,\ldots\,],
     editorial additions bracketed.}
  \end{center}
  \vspace{1em}}

\endinput


\notebook{black-notebook}
\ledgertitle{Josephine Hopper --- [Black two-ring notebook]}
\sheets{1}{68}
\dating{circa 1952--1961}
\watermark{Demonstration edition}

\begin{document}

\section{Page 1}

\sheet{96076}{1}
\hand{jo}{Mrs. Edward Hopper \\
3 Washington Square North \\
New York 3 \quad N.Y. \\
Phone \quad Spring 7 - 0949. \\
also Truro, Cape Cod, Mass- \\
\uncertain{Her book}}

\note{a box is drawn in ink round these six lines}

\hand{jo}{Phonograph \quad spring 1960 \\
\uncertain{Schirmers} \\
Arvin 7096 model. \\
Parts Code no. \quad 11 : \uncertain{46000} \\
Parts - \quad : 46229}

\note{a second drawn box encloses these five lines. The second figure of
« 7096 » is a barred nought, and the two part numbers are the least certain
figures on the page: what follows « 46 » in the first is three small loops of
one size}

\hand{jo}{\# 067-30-7265 \\
Jan. 20. 1954,}

\note{these two lines are in pencil, below both boxes and outside them}

\section{Pages 2--3}

\sheet{96248}{2}
\hand{jo}{Contes Comme ça - 95 \\
Chat qui marche + enfant d'élephant - \quad Troisième \\
Jan 16, 60.}

\note{« d'élephant » turns the line after « d'él- », with the hyphen written;
it is given whole here. Nothing else is on page 2}

\hand{jo}{Little white frames - inside \\
4 3/4" x 6 1/4" \\
4 of them on top shelf \quad rear studio \#3. \\
Put there Nov. 1948 - from \\
trunk in basement (\uncertain{burglars}) \\
Fitzgerald frames. \quad 51 1/2" x 35" \\
Wood panels - 10" x 14".}

\note{« rear studio \#3. » is written above the line, at the end of the third
line. The parenthesis is the least settled word of the seven: « burglars »
suits the strokes and the burglary of February 1958 is recorded on page 19,
five years later, which is not evidence for this line and does not settle it}

\hand{jo}{over 2 pages}

\hand{jo}{\ill{}}

\note{« over 2 pages » is ringed. Below it stand five arrows in ink and five
strokes of red crayon, and the lines they were drawn against have been rubbed
out; what shows under them is the show-through of page 4, which reads
backwards. Nothing in the rubbed block can be read}

\hand{jo}{cuts off bottom of canvas. \quad \uncertain{asking} E. \\
from back of frame \quad \ill{} \\
Silverfish - from Rehn Gal. on hill + pond \uncertain{backs} ? \\
Inside \quad 23 3/8 x \uncertain{37 1/2} \\
Shows - \quad 22 5/8 x 36 5/8 \\
what shows of canvas = 21 3/4 x 37 5/8.}

\note{« asking E. » is interlined above « bottom of canvas », and the word
after « from back of frame » is a single word of six or seven letters that
will not be read. A brace gathers the « Inside » and « Shows » lines and an
arrow runs from it to « what shows of canvas ». The numerator of « 23 3/8 » is
written over an earlier figure; the 3 of « 37 1/2 » carries the closed upper
loop that also reads as an 8, which the « 37 » two lines below does not}

\hand{jo}{20 x 24 - French Elaborate - Robinson 1948 \\
on San Estoban \\
29 x 36 fits in - but 5/8" covered all around \\
1950 - insert \uncertain{removed} to fit 29 x 36.}

\marginal{light}

\marginal{on Petunias on Cape.}

\marginal{Barbazon frame}

\note{« light » is interlined above « French », « on Petunias on Cape. » above
« Robinson 1948 », and « Barbazon frame » below the line in a darker ink. A
long horizontal stroke crosses « French »; it may be the t-bar of the
interlined « light » and it is not given as a cancellation here. « covered all
around » turns down the fore-edge corner. « removed » is followed by two
blotted letters. An arrow ties « on San Estoban » to the line above it}

\section{Pages 4--5}

\sheet{97063}{4}
\hand{jo}{To Ellen Ravenscroft \\
The Studio Gallery \quad 418. \\
2 W. 15" - 96" 5" Ave \\
Sat. Jan. 22, 44. \\
\emph{Water Colors}}

\hand{jo}{Provincetown from Websters \\
Shack + field - Fall's Village - gets new mat \\
Children in Gramercy \quad baby carriage \\
largish water color paper \quad 1/2 sheet \\
Boats - Gloucester high tide \quad Returning \\
\quad 2 in one mat \quad gold framed. \\
Boats - Gloucester \quad " \quad " \quad " \\
\quad - dories at \uncertain{barges} \\
\quad 5 water colors + 4 mats - \\
wrapped in brown paper.}

\marginal{travels}

\note{a brace gathers the « Children in Gramercy » and « largish water color
paper » lines. « travels » is interlined above « 4 mats ». The three ditto
marks on the second « Boats - Gloucester » line stand under « 2 in one mat »
and « gold framed » and are given as she set them}

\hand{jo}{Returned May 44 - \\
2 \quad Children in Gramercy - baby carriage \\
\quad High Tide Gloucester boats. \\
3 \struck{others} to be taken to Provincetown \\
June 1947 - these 3 not yet returned. \\
" \quad 1948 \quad Returned. \\
Ellen Ravenscroft died Jan. 1949, \\
1948 - left Lollipops water color at \uncertain{Maries} Sterners \\
not found Spring 1950. \\
In lot sold to Boston men after her death}

\marginal{Fall}

\note{a brace gathers the two lines after « Returned May 44 ». « Fall » is
interlined above « left ». The « " » opening the sixth line is a ditto for
« June »}

\hand{jo}{Du fond de son coeur - \\
\quad from bottom of one's heart \\
De tout son coeur. \\
\quad with all one's - \\
\quad de tout mon coeur. \\
one's = son \\
my = mon}

\note{« Du », « De » and « tout » are doubly underlined and « mon » once; a
brace gathers the last two lines}

\hand{jo}{Fountain pen - Altman - Spring 1956 - cost ? \\
\uncertain{Sheafer} medium flexible \\
L. E. Waterman \\
1 \uncertain{Detroit} St., Seymour, Conn.}

\note{a box is drawn round these four lines, and « cost ? » closes the first
against it}

\section{Pages 6--7}

\sheet{96080}{6}
\hand{jo}{\emph{To finish plain wood frames.} \\
\quad \uncertain{c̄} B. Hartman \\
1. Shellac. \\
2. White shoe polish liquid \\
3. Rub in dry color - pastel - \\
\quad \uncertain{remove do.} \\
4. Wax.}

\note{the mark before « B. Hartman » is a small letter under a macron, the same
mark that stands in « Truro, trip c̄ Manchester » on page 19. A short rule
closes the four numbered lines}

\hand{jo}{\emph{Water Colors -} \\
SanEsteban Interior \quad 30 x 22 3/4 \\
Siera Madre + roofs from Monterrey \uncertain{Hotel}. \\
\quad 24 + or - x 30 \\
Wittenberg Garden - 20 x 24 \\
Ascension \quad 20 1/2 x 25 1/2 \\
Shelf paper - 13 x 19 ? \quad Mats 19 x 25 \\
Watman" 1/2 sheet - 15 x 22.}

\note{« Hotel » is written over an earlier word and is not certain. The place
is « SanEsteban » here, written as one word, and « San Estoban » on pages 3 and
10}

\hand{jo}{Oil - Bedroom Wash Sq., desk, clock, fireplace - \\
\quad 27 x 32 \\
Big stove - 32 x 38 - Picture size. \\
Small oils - 10 x 14.}

\section{Pages 8--9}

\sheet{97564}{8}
\hand{jo}{Water color paper - \quad 22 x 30 \\
\quad 1/2 sheet = \quad 15 x 22. \\
Water color pad - \quad 14 x 20. \\
Lindenmeyer mat \\
paper - - - - - \quad 25 x 38. \\
\quad 1/2 sheet - - - \quad 19 x 25 \\
Grey pulp board 1/2 sheet \quad 20 1/2 x 31 \\
\quad 1/2 \quad " \quad " \\
1/2 Lindenmeyer \struck{mat} paper \quad 19 x 25 \quad mat.}

\note{« paper » is interlined above the struck « mat » in the last line, and
the 19 of « 19 x 25 » is written over an earlier figure}

\hand{jo}{1/2 water color \quad " \quad 15 x 22 \\
\quad 2 )4 \qquad 3 \\
\quad Margin - = \quad 2 \qquad 1 1/2 \\
\quad to small. \\
Black portfolios \quad 20 x 26 1/2}

\note{« 2 )4 » is set as a division, the 4 under a rule}

\hand{jo}{Oil sketch box panels - old \quad Early - \\
Prov. from Wharf. \quad 10 x 14. \\
Sail boats tangled spars \quad 10 x 14. \\
From Websters. Provincetown - Bay. \quad " \\
Quincy Boats - \qquad " \\
7 framed ones - Big Sail, Red Masts, Dories on White beach \\
Clothes Line, Little Sail boat touching beach, White sail coming}

\marginal{part}

\marginal{yellow \quad " \quad " \quad "}

\note{« part » stands above « Dories on White beach » and « yellow » with three
dittos above the last line. A short addition runs up the fore-edge beside these
lines, of which only « w.c. Gloucesters » can be made out and the rest is
crowded against the edge and unread}

\hand{jo}{\emph{Frames.} \\
Usui - May, 53 - 25 x 30. \\
Chez Hopper \quad - \quad 29 x 36 \\
Fruit - \qquad 20 x 24 \\
Church - \qquad 22 x 38. \\
\uncertain{Dauphines} House \quad 20 x 28. \\
Sommerfield frame \quad 22 x 38. \\
\uncertain{Saidelyn} (Judson ones) \quad 29 x 36 \\
Robinson - fancy French. \quad 20 x 24 \\
Heidenryk - \uncertain{epoxy} white - less than}

\marginal{Early}

\marginal{29 x 36 1/2 \quad 27 x 34}

\note{« Early » is interlined above « May » in the Usui line. The two
dimensions of the second marginal are written up the fore-edge beside the
Robinson and Heidenryk lines and belong to them. The frame maker is
« Heidenryk » here, « Heinrich » on page 10 and « Heydenryk » on page 18}

\hand{jo}{Water color paper. \\
\quad Wat. mat's \quad roughs \\
\qquad 22 x 30" \\
Coper Union - 72 lbs. - \quad 20 ¢ \\
\qquad 140 " \qquad 39 ¢ \\
See Robinson - \quad 72 " \qquad 23 ¢ \\
\qquad 90 " \qquad 37 ¢ \\
\qquad 140 " \qquad 50 ¢}

\note{a short addition runs up the fore-edge beside the « Wat. mat's » line, of
which only « a sheet » can be read; the rest is crowded against the edge and
unread}

\hand{jo}{\emph{Mrs Fitzgerald's frame.} \\
\quad 30 3/4 x 41 1/2. \\
Old Silver on \emph{Back of Cape Houses} - pond, sandhills, church \\
\qquad 37 1/2 x \uncertain{23}}

\note{the last figure carries a small raised mark that may be the numerator of
a fraction whose denominator is not on the page}

\section{Pages 10--11}

\sheet{97318}{10}
\hand{jo}{Barbizon frames \quad 20 x 24 - 3. \\
20 x 24 \quad Cape Bedroom \\
\qquad Obituary - dead flowers, cat, \\
\qquad Petunias - on Cape \\
29 x 36 - made for \uncertain{St. Frans Ch. Buxton} \\
27 3/4 x 34 3/4 - 1" white insert \uncertain{of}. \\
on San Estoban - \quad Heinrich grey white frame.}

\marginal{grey - antique.}

\marginal{Robinsons}

\marginal{1949}

\marginal{sea}

\note{a brace gathers the three lines under « 20 x 24 ». « grey - antique. » is
interlined above « Barbizon frames » and « Robinsons » stands at the head of
the page above it; « 1949 » is above « made for »; « sea » above « dead
flowers ». She writes « Barbizon » here and « Barbazon » on page 3}

\hand{jo}{Mat on San Estoban Interior - (water color) \\
\quad 30 x 22 3/4.}

\note{a box is drawn round these two lines}

\hand{jo}{Need frame for Bed room, desk, fire place, clock \\
\qquad 27 x 32 \\
Mexico - Sierra Madre + roofs - San Juan \\
\quad 24 (+ +) x 39 could be cut \\
2 gold from Rehn Gal. \quad 28 1/4 x 36 1/2 \\
Silver frame from Rehn Gal. on Hills + Pond back of house \\
\quad = 37 1/2 x \uncertain{23 1/4}. \\
Canvas 36 x more than 23 1/2 \\
Hartman frame - outside = 17 1/2 x 24 1/2 \\
\quad not large enough. \quad Matt frame \\
Frame on Ascension \qquad " \qquad "}

\marginal{phone.}

\marginal{wider}

\note{« phone. » is interlined above « fire place »; « wider » above « could be
cut ». A tick stands to the left of the « Mexico » line and of the « 24 (+ +) »
line. She writes « Sierra Madre » here and « Siera Madre » on page 6}

\hand{jo}{\emph{Insurance \quad 1947, 48} \\
So. Truro - Family Theft \quad F.T. 62458. \\
Amer. Employers, Edward Stone Pres. \quad Boston Mass. \\
Expires Dec. 18, 1950. \qquad taken out Dec. 18, 47. \\
Premium - 22.50 \\
\uncertain{Series} \$1000 \quad Prem. 22.50. \\
Standard Windstorm \quad \# \uncertain{T} 3387 \\
N. British \& Mercantile Ins. Co \\
Cottage + Garage - \qquad 1000. \\
\qquad Prem. \qquad 5. \\
Expires Jan 30, 51.}

\marginal{\struck{yearly}}

\note{the whole of this side is cancelled with a large cross drawn over it in
pencil, and the two « Expires » lines are each ringed in pencil. « yearly » is
interlined above « Premium - » and struck. The first word of the « \$1000 »
line is the least settled on the page}

\hand{jo}{N.Y. League for the Hard of Hearing \\
T 4 L \quad test \\
(solution)}

\marginal{for.}

\note{these three lines are in pencil, « T 4 L » in large block capitals, and
they fall below the cancelling cross. « for. » is interlined above « test »}

\section{Pages 12--13}

\sheet{97145}{12}
\hand{jo}{\emph{Theatres with 2" balcony.} \\
Henry Miller \quad E. of 6" \quad 43" St. \\
Empire - Bd. \quad 39" \\
Belasco \quad 44" E. of 6" \quad Gentle People \\
Fulton \quad W. 46" \quad Oscar Wilde \\
Playhouse. W. 48" - Outward Bound \\
Windsor \quad 48 E. of 6" \quad Awake + Sing. \\
Court - W. 48" \qquad Primrose Path \\
Lyceum - 45" N. of 6" \quad When we are married}

\marginal{nice}

\marginal{arsenic + \uncertain{gold lace}}

\note{« nice » stands in the left margin against the Fulton line and « arsenic
+ gold lace » above it; the second word of that interlineation is written as
one and begins with a clear g. « 2" » and the street numbers carry her raised
ordinal, which is the same mark she uses as a ditto}

\hand{jo}{\emph{Avoid} \quad St. James \quad W. 44" \quad "12" Night" \quad 2" balcony far \\
\qquad back of 1" balcony, \\
(\emph{Avoid} the Biltmore - W. 47" St \\
long, narrow balcony - 1 balcony, \\
\qquad neither see nor hear) \\
("Flashing Stream" - wasted.) \\
48" St. theatre. Excellent. \quad "row 2" (Harvey) \\
\emph{Avoid} Civic Centre \quad 2" balcony - high up. \\
\qquad also back of 1" balcony \\
Court - 48 E. of Bd. \quad fine 2" Gal. \quad Season in the Sun. \\
\qquad Mar. 17, 51 \\
\emph{Avoid} \quad Ethel Barrymore - \\
\qquad 7 Balcony is one block from stage \\
\qquad Four Poster - Nov. 51 \quad no 2" bal. \\
\emph{Avoid} \quad Ziegfeld - Caesar + Cleopa. \quad Oliviers \\
\quad 2 blocks deep - 1 balcony + Mezzanine \quad Jan 1952 \\
\quad deeper than " \quad " \qquad " \quad "}

\note{each « Avoid » is doubly underlined. Blue crayon underlines « Civic
Centre », « Court » and « fine 2" Gal. »; the crayon is hers and is not given
in the text. « 12" Night » writes Twelfth with her ordinal mark. The mark
opening the « Balcony is one block » line is a 7 or a long hooked dash and is
given as a 7}

\section{Pages 14--15}

\sheet{96257}{14}
\entry{Spring 1946.}{1946}
\hand{jo}{Sam'l \\
\emph{Golden Amer. Artists Group Monographs} \\
\quad given to; \\
\uncertain{LeB. Bakers} \\
Louise Floyd. \\
Jean Chamberlin \\
Dr. Thos. E. \uncertain{Dwyer} \\
Edna Stauffer \\
Anne Tucker \\
Bee Blanchard \\
Howard + Anne \ill{} \\
Horace + \uncertain{Geo.} \quad " \\
Marion Hopper \\
\quad J \& N. \quad " \\
Theodore \uncertain{Sicanteri} \\
Snta. Dolores Gomez - Sortillo \\
Ella Mielziner - Xmas. 46. \\
Dan Stephens, Bob + Marie \\
Frantz Lombard 1947 \\
Ed. D. + Louise Moore \quad Winter, 47 \\
Dr. Earl \uncertain{O'Keil} - N.Y. Hos., Feb. 48 \\
Nurse \uncertain{Hélène Felenfeld} \quad " \qquad " \\
" \quad Augusta Brand. \qquad Mari \\
Nurse Barton \\
Miss Harrington \quad for her own \\
Catherine Rogers - June 48}

\hand{jo}{Hugh + Dr. Ferris. \quad June 48 \\
Edna Houghton \\
Dick + Nell Magee \quad Aug. 48 \\
Redmond Lane \quad Fall 47 \\
\struck{Magee's friends - 48} \\
Enid Buhre \quad Sweden \\
Frances Manning \\
Frances Baldwins \\
Dorothy Perry \uncertain{Juchin} \quad 49 \\
Jane Wilson \uncertain{(Sfo)} \quad Feb. 20 1950 \\
Angel Barnes \quad Insurance \\
\quad Mexico \quad 1946 \\
\quad \ill{} \quad \uncertain{Thomason}}

\marginal{Marion friends}

\marginal{Summer}

\marginal{Sailors Snug}

\marginal{English}

\marginal{27 \qquad 29 \qquad 46}

\note{the page is two columns and the second was written later, smaller, and
partly across the first, so the pairing of a name in one column with a name in
the other is not always recoverable and is not asserted here: the two columns
are given in order, left then right. The surname after « Howard + Anne » is
overwritten by « Frances Baldwins » and cannot be read. « Mari » breaks off,
and a long hooked stroke runs from it to « over » at the fore-edge. Loose
numerals - 27, 29, 46 - stand between the columns, and « 11 » in the left
margin beside « J \& N. ». « Magee's friends - 48 » is struck through with a
long stroke. « Nurse Barton » has a rubbed pencil word above it. The last
three lines of the right-hand column are the most crowded on the page and
« 1895 4C Thed. Co » or some such stands unread between « Mexico » and
« Thomason »}

\section{Pages 16--17}

\sheet{97045}{16}
\hand{jo}{\emph{E. Suit Size}}

\begin{ledgertable}{Y{0.34}rr}{ & & 1952}
Waist & 38 & 40 - \\
Chest & 42 + & \\
Leg & 35 - & 34 \\
Length of Coat & & 34 \\
\end{ledgertable}

\note{she ruled the two right-hand columns herself and headed only the outer
one, « 1952 »; the two left headings are hers to leave blank and are not
supplied. The whole block, heading and table, is in pencil inside a drawn box}

\hand{jo}{Neck - 15 1/2. \quad but work shirts 16.}

\hand{jo}{\emph{Monographs to} \\
Harry \qquad Dauphinée - Nov. 8, 49.}

\note{« Dauphinée » here, « Dauphines House » on page 9}

\hand{jo}{the Jarvis Commonweal Fund \quad (Jean Chamblin) \\
158 Bowery \quad (Miss Dennis) \\
People over 65 with cultured background \\
too good for old ladies' homes \uncertain{subsidised}. \\
Miss Jean Chamblin - ± 1933 or earlier \\
" \quad Lillian Childs - later \\
\quad Benficiaries. We sponsored J.C. \\
the Jarvis Commonweal Fund. \\
Refound Oct. 14", 1964.}

\note{a brace gathers the two « Miss » lines. « Benficiaries » is as written.
« ± » is her mark for about}

\hand{jo}{\emph{Old frames} \\
Old Whistler moulding - in rack - \\
\quad 38 x 48. \quad heavy with insert \\
\struck{Frame of \uncertain{Peters fath.} - the \uncertain{Cole}} \\
\quad 29 x 36. \\
(Chez Hopper \quad " \quad size.) \\
26 x 38 1/2 \quad insert of 1 1/2" - so \\
\quad could be used on larger canvas \\
24 x 29 \qquad \uncertain{wid.} insert \\
32 x 38 1/2 \qquad " \qquad " \quad Big stove = 32 + 38. \\
20 1/2 x 29 \qquad " \qquad " \quad narrow - not so good \\
24 x 29. \\
\quad 32 1/4 x 38 1/2 \\
24 x 29 - could use on 25 x 30 \\
\qquad restretched ? \\
\qquad Cat Madonna ?}

\marginal{including}

\note{« 29 x 36. » is written across the struck line, not under it. A tick
stands to the left of the « 26 x 38 1/2 », « 24 x 29 », « 32 x 38 1/2 »,
« 20 1/2 x 29 », « 24 x 29. » and last « 24 x 29 » lines, and « including » is
interlined above « insert of 1 1/2" ». A brace gathers the « 20 1/2 x 29 »
line, and a long hooked arrow runs from the left margin down to the last
« 24 x 29 »}

\section{Pages 18--19}

\sheet{97168}{18}
\hand{jo}{\emph{New frames} \quad 1952 \quad Matt. \\
Water Color - on \quad made by Matt - E. 54" St. \\
Ascension - 20 x 25 \quad \uncertain{+} mat \\
Bertram H. \quad 14 x 20 1/2 \qquad 18 x 25 \\
19 1/2 x 26 1/2 \qquad Mat 2" - bottom 2 1/2 \\
outside \qquad 18 x 25 inside \quad frame 1 1/2 \\
Silver frame f/ Canvas from Rehn Gal. \\
\quad 37 1/2 x 23 1/4 \qquad Canvas = 36 x 25 \\
Mrs. Fitzgerald - Nyack - old frame at Marion's \\
\quad 30 3/4 x 47 1/2 \\
Usui frame 25 x 30 - gold, \uncertain{pinkish} - \$18 \\
Heydenryk - 25 x 30 - goldish + white. \quad 11.50 \\
Matt - 25 x 30 was aluminum + gold + no good \\
on anything. Too raw sienna - awful in \\
north light. Gorgeous glittery in fluorescent \\
now worse - \emph{more} yellow than before.}

\marginal{\$14}

\marginal{\$16}

\marginal{inside frame.}

\note{« \$14 » is interlined above « Ascension » and « \$16 » above
« Bertram H. »; « inside frame. » stands above « 18 x 25 » in the fourth line.
Two braces gather the Bertram H. block. The mark before « mat » in the third
line is a cross with a second stroke through it and is not her plain +. The
cents of « 11.50 » are written raised. A red crayon stroke underlines the
« 37 1/2 x 23 1/4 » line and another stands in the left margin beside it}

\hand{jo}{Ravel - Cape Bedroom - 29 x 36 \\
Wash. Sq. Bedroom - June 1956. 27 x 32 \\
Cape Petunias \quad 20 x 24 \\
Portrait - E. H. not 12 x 16 - 11+ x 15+ \\
7 \uncertain{Pochardes} - 10 x 14. \\
Treetops beyond fire escape - 27 x 43 \quad 1960 with gold.}

\marginal{insert}

\note{a brace gathers the four lines from « Wash. Sq. Bedroom ». « insert » is
interlined below « 12 x 16 »}

\hand{jo}{\emph{Frames:} \\
Pochards - 10 x 14 - 1956 June 25" \quad \uncertain{Each priced} \\
Ravel. No. 214 \quad (1) \quad 18 x 14" finish \quad 19 \quad \uncertain{12.05} \\
\qquad 241 \quad (5) \qquad 819 beige \quad 8.30 \\
Order \# 08722 - June 25, 1955}

\marginal{\struck{It frame}}

\note{the cents of « 12.05 » and « 8.30 » are written raised, and the last two
figures of « 12.05 » are written over others}

\hand{jo}{\uncertain{Path} from Studio \quad 200 - 29 x 36 finish. 2R = 30.75 \\
Fire Escape \qquad 200 \quad 25 x 30 \quad 0 = 26.50 \\
Feb. 3, 58 \quad Order No. 04413}

\marginal{frame no}

\marginal{\uncertain{whimsical}}

\marginal{\uncertain{wh. mix}}

\note{the cents of « 30.75 » and « 26.50 » are written raised. The two
interlineations above the « Path from Studio » line and the one above « 25 x
30 » are small and doubtful}

\hand{jo}{Rack frame in dumbwaiter closet \\
\struck{\ill{}} \qquad 46 x 35 1/2" minus \\
46 x 35 1/2" \qquad Dec. 31, 1960.}

\note{the struck matter between « closet » and « 46 x 35 1/2" » is a dimension
written over and then crossed, and cannot be read. The lower « 46 x 35 1/2" »
is boxed in red crayon}

\entry{1961.}{1961}
\hand{jo}{Boston Mus. T.V. \quad Ap. 20 \quad (N.Y. \quad Mar. 29)}

\entry{1960.}{1960}
\hand{jo}{S. Truro - \quad \uncertain{(Sr. Boston Mar. 1, 2)}}

\entry{1959.}{1959}
\hand{jo}{Truro, trip c̄ Manchester. \ill{} \\
N.Y. Hos. return Jan. 23, 58}

\marginal{a Providence - Dec. 2 nights. N. Chace}

\marginal{2 nights - Oct. 59}

\marginal{Eliot 3 Cent'y Aud.}

\entry{1958}{1958}
\hand{jo}{- burglar Feb. 11", 58 - \\
Feb. 17+ Truro. Dec. 17. N.Y. Hos. Jan. 29, 58}

\marginal{Sun. Assn. Life Nov. 17}

\entry{1957}{1957}
\hand{jo}{July - 17 ± Truro. Dec. 17. N.Y. Hos. Jan. 29, 58 \\
Dec. 05 1956 - June 6, 57 at \uncertain{H. H.} + Pacific Palisades}

\marginal{Cal.}

\note{a brace gathers the two 1957 lines. The lower half of page 19 is in
pencil and is crowded line against line; the marginals above are the
interlineations of the years they follow, and the placing of one or two of
them against a particular year is not certain}

\section{Pages 20--21}

\sheet{97496}{20}
\hand{jo}{Dec. 6" 1956 \quad H. H. + Pacific Palisades}

\note{this line repeats page 19's last, which is hers and is not merged}

\entry{1956}{1956}
\hand{jo}{- Coke, Mex. \\
Monterrey - lv. May 31 \quad return May.}

\marginal{Nov. start 11 for S.W. - stop Tuscon, Ap. El Paso}

\marginal{Motel Anza \quad Ap. 9 - May 1 to Mx.}

\marginal{in new Buick}

\entry{1955 -}{1955}
\hand{jo}{Truro - July - Nov.}

\entry{1954 -}{1954}
\hand{jo}{S. Truro. E. N.Y. Hos. hernias - Feb. \\
to S. Royalton \quad Slaters \quad Charlestown \\
July - S. Truro \quad Sept. Gloucester trees \quad Frank Rehn}

\marginal{brief visit}

\marginal{staying \quad Elsa on Cape}

\marginal{Woodstock Vt. - Katherine Proctor}

\entry{1953}{1953}
\hand{jo}{Mexico back May. - Richmond for Va. jury \\
New Brunswick - Dr. Sutters. Rutgers. \uncertain{Butler Mus. O.} \\
at Youngstown, O. \quad June \qquad Nyack. \\
Mexico \quad Lv. Dec. 1", 52 - El Paso, Tx. \quad Oaxaca}

\marginal{May 8}

\marginal{Youngstown O.}

\entry{1952 -}{1952}
\hand{jo}{Cape. By air to N.Y. for Metrop. jury \quad Sept. 52 \\
\quad Saltillo \quad Electric - Silversmiths \\
Mexico + Cape, Truro back in July}

\marginal{Efforts}

\marginal{to light, electric}

\marginal{May 28. Santa Fe}

\entry{1951.}{1951}
\hand{jo}{Washington Biennial Corcoran \quad Ap. 25 \quad March 2" \\
Truro \quad Trips to Boston, Chicago, Detroit \\
E. gets degree Dr. of Fine Arts - Chicago Art Inst. \\
Big Retrospective at Whitney,}

\entry{1950 -}{1950}
\hand{jo}{3" trip N.Y. Hos. \quad Diverticulisis \\
\quad S. Truro.}

\entry{1949 -}{1949}
\hand{jo}{2" trip N.Y. Hos. \\
\quad S. Truro -}

\entry{1948 -}{1948}
\hand{jo}{1" trip to N.Y. Hos. \quad \uncertain{We} Show at Rehn - \\
\quad July to Truro.}

\marginal{8 canvases}

\entry{1947 -}{1947}
\hand{jo}{\# 3 invasion by N.Y.U.}

\entry{1946.}{1946}
\hand{jo}{Sattillo, \quad Sq. G. Camp \quad \uncertain{Bravos Colons}. \quad S. Truro.}

\note{page 20 is the densest of the twelve sheets and carries an interlineation
above almost every line; the marginals are given under the year whose lines
they stand over. She spells the town « Saltillo » in the 1952 entry and
« Sattillo » in the 1946 one, and « Sortillo » on pages 15 and 21. Tucson is
« Tuscon »}

\entry{1945 -}{1945}
\hand{jo}{Cape - ap. Salmagundi Prize. \\
Xmas 44 \quad E. elected to Instit. of Arts + Letters}

\marginal{mar. - \uncertain{Du Barry caps.}}

\entry{1944 -}{1944}
\hand{jo}{Cape - Boston trip in car.}

\entry{1943 -}{1943}
\hand{jo}{Spring - Washington - Corcoran \\
\quad Summer - Mexico City \\
\qquad Sortillo \\
\qquad Monterrey}

\marginal{jury}

\note{a brace gathers the three place names}

\entry{1942 -}{1942}
\hand{jo}{Cape.}

\note{a line of pencil under this entry has been rubbed out; a few strokes and
what may be « trip ? » show, and nothing in it can be read}

\entry{1941 -}{1941}
\hand{jo}{Trip West - Calif., Oregon, Wyoming. \\
\quad back to Cape}

\entry{1940 -}{1940}
\hand{jo}{Cape \quad back to register. \\
\quad late fall - Albany trip \\
\quad or spring 41.}

\entry{1939}{1939}
\hand{jo}{Cape - left early for \\
\quad Pittsburgh trip. Carnegie jury}

\entry{1938 -}{1938}
\hand{jo}{Cape \qquad Huricane \\
\quad Vermont - Slaters \quad 2" visit. \\
spring - Richmond ?}

\marginal{back to Cape.}

\entry{1937.}{1937}
\hand{jo}{Cape \\
\quad Spring Washington. \quad jury at Corcoran. \\
\quad Vermont Slaters}

\entry{1936 -}{1936}
\hand{jo}{Cape. \\
\quad Vermont ?}

\note{« Huricane » is as written}

\section{Pages 22--23}

\sheet{97026}{22}
\entry{1935 -}{1935}
\hand{jo}{Cape.}

\marginal{Harold + Florence Weed}

\marginal{+ later - E. Montpelier, Vt.}

\marginal{+ back to Cape.}

\note{the first of these three stands to the right of the 1935 line and the
other two below it, the last running under the red rule at the head of the leaf}

\entry{1934 -}{1934}
\hand{jo}{Cape - Built house.}

\entry{1933}{1933}
\hand{jo}{Quebec + stop at all old haunts \\
along Coast - on to \\
Chestnut Hill over night.}

\marginal{Pregnant Buick new tires}

\entry{\uncertain{1933}}{1933}
\hand{jo}{Cape - last yr. at Bird's Cage. \\
Bgt. land.}

\marginal{Mus. Mod. Art Show - E. Retrospective.}

\note{the year opening this entry is rubbed and stands faint; it reads 1933 a
second time and is not the stroke of the 1933 above it}

\entry{1932 -}{1932}
\hand{jo}{Cape - Bird's Cage}

\entry{1931 -}{1931}
\hand{jo}{Cape - Bird's Cage}

\entry{1930 -}{1930}
\hand{jo}{Clinton - tourist Roots \\
Cape - Wellfleet then \\
Bird's Cage S. Truro.}

\note{« Roots » is as written}

\entry{1929.}{1929}
\hand{jo}{Charleston, S.C. \\
Topsfield, Mass. - S.A Tuckers \\
trip N. - to Jones Point. - \\
Two Lights. - at Maxwells.}

\entry{1928}{1928}
\hand{jo}{Gloucester - Middle St.}

\entry{1927 -}{1927}
\hand{jo}{Two Lights - Coast Guard Cottage \\
Movies at \uncertain{Pattesons} studio \\
got car - Dodge \qquad Ogunquit - Charleston}

\marginal{3" aniv.}

\entry{1926}{1926}
\hand{jo}{Gloucester - Middle St. \\
\quad after Eastport, Me.}

\entry{1925 -}{1925}
\hand{jo}{Santo'Fe' + New Mexico \\
\quad James Mt.}

\entry{1924}{1924}
\hand{jo}{Gloucester. Honeymoon. \\
Watercolors. Chicago August.}

\section{Pages 24--25}

\sheet{96128}{24}
\marginal{\uncertain{Grece} House sold \\ Bklyn. Mus.}

\note{this stands at the head of the leaf, above the 1923 line, and the first
word will not settle between « Grece » and « Green ». A monogram is drawn in a
circle in the top corner beside it}

\entry{1923 -}{1923}
\hand{jo}{Gloucester +}

\entry{1922 -}{1922}
\hand{jo}{Provincetown}

\marginal{\uncertain{Sept.} N.J.}

\entry{1921}{1921}
\hand{jo}{Woodstock}

\marginal{1920 - \uncertain{Vandenberg} Studios}

\marginal{1919 - 50 Chapel 3 Univ. Pl.}

\note{these two are written one above the other across the end of the 1921
line, and carry years of their own that do not agree with the years they sit
against}

\entry{1920}{1920}
\hand{jo}{Aloha Hive, Lake Fairlee, Vt.}

\entry{1919}{1919}
\hand{jo}{Ogunquit \qquad Bret + Washington D.C. \\
\qquad Bean Desert, Nov. 18.}

\note{a brace gathers the two right-hand phrases of this entry}

\entry{1918 -}{1918}
\hand{jo}{Monhegan -}

\entry{1917 -}{1917}
\hand{jo}{\qquad Provincetown ? Fish House?}

\entry{1916}{1916}
\hand{jo}{\qquad Ogunquit. ?}

\entry{1915 -}{1915}
\hand{jo}{Monhegan ? \qquad Wash Sq. Players.}

\note{a brace stands between « Monhegan ? » and « Wash Sq. Players. »}

\entry{1914 -}{1914}
\hand{jo}{Provincetown + Ogunquit}

\marginal{War broke out}

\entry{1913}{1913}
\hand{jo}{}

\entry{12}{1912}
\hand{jo}{}

\entry{11}{1911}
\hand{jo}{}

\entry{10}{1910}
\hand{jo}{Martha's Vinyard + Siasconset}

\entry{9}{1909}
\hand{jo}{Provincetown - no Ariel.}

\entry{8}{1908}
\hand{jo}{Provincetown - Ariel.}

\entry{7}{1907}
\hand{jo}{Haarlem, Italy}

\entry{6 -}{1906}
\hand{jo}{Provincetown}

\note{the chronicle begun on page 19 ends here. The years 1913, 12 and 11 stand
with nothing written against them, and are left blank rather than passed over}

\hand{jo}{Chrysoprase = green stone in Greta \\
44 coral beads , Silver brooch.}

\marginal{Kuehne's}

\note{« Kuehne's » is interlined above « beads », and belongs to « Greta » at
the end of the line above. A pencil rule is drawn under the two lines}

\hand{jo}{For burns - tannic acid in tube.}

\hand{jo}{Who's Who in Amer. Art - 1947 \quad oo \quad (46?) \\
Amer. Fed. of Arts \qquad 1953 \\
Rev'd for them by R.R. Bowker Co. \\
\quad 62 W. 45", N.Y. 36. \qquad (Litt. incl) \\
\quad Dorothy B. Gilbert. Editor}

\note{« (46?) » is circled, and so is « Litt. incl »}

\hand{jo}{Jean C. - \qquad Jarvis Foundation (?)}

\note{this line is boxed in pencil and underlined in red crayon; the question
mark that closes it is hers and is circled}

\hand{jo}{Ret. Oct. 1, 21. \quad 14 - 5 - 5. \quad '85. \\
\quad Feb 3, 53.}

\note{two vertical strokes divide the first line into three, before « 14 » and
before « '85 »}

\entry{« Reality » - Ap. 17 ±, 1953}{1953-04-17}
\hand{jo}{May 3, 53 - J. Donald Adams.}

\marginal{N.Y. Times Book Review}

\marginal{Speaking of Books}

\note{a red crayon stroke stands in the margin beside « Reality ». The mark
after « 17 » is a plus over a bar and is hers}

\hand{jo}{Studio + Bedroom painted 2 Coats - Casein \\
\quad May 19 1953. - 4 day job. \qquad \$250 or 350?}

\marginal{+ bath}

\marginal{water paint}

\note{« + bath » is interlined above « painted » and « water paint » above
« Casein »}

\hand{jo}{Rutgers - D. Lit. June 3, 53. honorary degree E.}

\section{Page 26}

\sheet{96131}{26}
\hand{jo}{\emph{Eyes - 1952}}

\begin{ledgertable}{Y{0.30}rrr}{Lines & R. & L. & Together}
Jan. 29" \quad \emph{10 ft.} & 3 + & 1 + & 5 - \\
\quad 31 & 4 & 1 - & 4 \\
Feb. 6" & 4 (-) & 1 - (double!) & 4 - \\
Feb. 12 & 3 - & 1 - & 3 + \\
Ap. 18, 52 & 6 & 3 - double & all 7, last not \\
\end{ledgertable}

\marginal{after painting all a.m. detailed work.}

\marginal{\emph{6 ft.}}

\marginal{tendency to}

\marginal{so easily}

\note{she ruled the three vertical lines of this table herself and wrote all
four headings. « after painting all a.m. detailed work. » is interlined between
the 31 row and the Feb. 6" row; « 6 ft. » stands above « Ap. 18, 52 »;
« tendency to » above « 3 - double » and « so easily » above « all 7, last
not ». The dash in the R. cell of the Feb. 6" row is circled. Page 27 is blank}

\section{Pages 28--29}

\sheet{96130}{28}
\hand{jo}{\emph{Photos at Times - Ap. 1953.} \\
Homage to Ylla - Cat Madonna. \\
Chez Hopper I. \\
Fire Place \\
Fire Escape + bird \\
Judson Tower \\
Obituary in frame \\
Park from window. \\
Coal Scuttle \\
Glacial formations \\
San Juan ? \\
Truro Hills + pond. \\
Sleeping Cat water color.}

\hand{jo}{\struck{Had} San Esteban \quad Summer 1957 to Editor}

\marginal{sent}

\note{« sent » is interlined above the struck « Had »}

\hand{jo}{\emph{Arts + Letters - 1955.}}

\begin{ledgertable}{Y{0.28}Y{0.34}Y{0.28}}{ & & }
Bertram & Chez Hopper I & Fanny \\
San Esteban & Fire Escape + bird & Obituary \\
San Juan & Truro Hills - glacial & Judson Tower \\
Clothes Line & " \quad " \quad + pond & Park from window \\
\end{ledgertable}

\marginal{must bottom chain}

\note{these three columns are not ruled; they are read downward, and the dittos
in the middle column stand for « Truro Hills » in the line above. No heading is
written on the leaf and none is supplied. « must bottom chain » is written
below « Park from window » and bracketed to it}

\hand{jo}{\emph{Reproduced in Villager} \\
San Esteban \\
Fire Escape + Convent opposite \\
Cape Clothes Line \\
Power of The Press - \uncertain{Harrocks} - 1958 \\
Park from studio window \qquad "}

\note{a brace gathers the five lines under the heading}

\hand{jo}{25 Wash. Sq. N. \\
\struck{244 E. 18" \qquad \uncertain{GR.} 3 - 0160.}}

\hand{jo}{\emph{Oliver Baker photos. \quad May 19, 50.} \\
10189 \quad Obituary + frame \\
\quad 90 \quad Park + window with chairs \\
\quad 91 \quad "Jewels for the Madonna \\
\quad 92 \quad Fire Escape + sparrow \quad 25 x 30 \\
\quad 93. \quad Fire place. bedroom \# 3 Wash. Sq. \\
\qquad 32" x 27" \\
6760 \qquad San Esteban + others \quad 29 x 36}

\marginal{20 x 24}

\marginal{29 x 36}

\marginal{roughly}

\marginal{Homage to Ylla}

\marginal{25 x 30+}

\marginal{\ill{} 1/5 \quad bottom}

\note{« 20 x 24 » stands above « Obituary + frame », « 29 x 36 » and
« roughly » above the 90 line, « Homage to Ylla » above the 91 line and
« 25 x 30+ » above the 92 line. A curve gathers the three upper entries and
carries the last marginal at the fore-edge, whose first figure is circled and
will not be read}

\hand{jo}{10245 \quad Watercolors) \quad Bertrams Hartman \\
10247 \qquad \uncertain{Quincy} Fleet \\
10246 \quad 20 x 24 \quad Ruth Wittenberg's garden \\
10248. \quad (oil) Glacial Formations}

\note{a long curve gathers these four entries at the fore-edge}

\hand{jo}{10244 \quad Calvin \uncertain{How} + his Sisters}

\marginal{May 28, 52}

\hand{jo}{\emph{July 9, 52} \\
10484 \quad San Juan. 38 x 28 ? \\
10,485 \quad Judson, bus + fire escape 29 x 36 \\
10486 \quad Overdressed Lily \\
10487 \quad Coal Scuttle. Chez Hopper 20 x 24 \\
10488 \quad Truro Hills, + Pond \quad 25 x 36 \\
10489 \quad Sleeping Cat - watercolor.}

\marginal{Church}

\marginal{gr. not}

\marginal{taken in 9}

\marginal{1st frame}

\note{« Church » is interlined in the 10488 line after « Pond » and « gr. not »
above its dimensions; the last two stand above and after « watercolor. » in the
10489 line, and are small}

\hand{jo}{(Not Baker - Chez Hopper. I (feet) \\
\quad " \quad ave.}

\marginal{old print}

\section{Pages 30--31}

\sheet{97590}{30}
\hand{jo}{\emph{Oliver Baker photos con.} \\
Rec'd. May 6, 53. \\
Chez Hopper I with feet. \qquad 12 - 558 \\
" \qquad " \quad Stove + umbrella \qquad 12 - 557}

\marginal{29 x 36}

\marginal{25 - 17}

\note{the two dimensions are interlined above the figures at the ends of the
two lines}

\hand{jo}{\emph{June 14", 54.} \\
15 - 133 - Cape Clothes Line \\
15 - 134 \quad Cape Bedroom - 2 windows \\
15 - 135 \quad N.E. corner of Sq. \\
15 - 136 \quad Fanny + periodicals \\
15 - 137 \quad Our Lady of Good Voyage}

\marginal{29 x 36}

\hand{jo}{\emph{June 27, 1957.} - \quad not as good as usual. \\
Camp Ground for Battle of Washington Sq. \quad 23.353. \\
x Hill close to studio window \qquad 23 354 \\
Village Car on Pacific Coast - \qquad 23 355 \\
x Mountain + Sanctuary \qquad 23 356 \\
x Petunias \qquad 23 357 \\
Betty Carmer - water color \qquad 23 358 \\
Head of \uncertain{Keith} \qquad 23 359}

\note{a word after « Head of » is rubbed out and cannot be read}

\hand{jo}{Jewels of the Madonna. \\
Fire Escape - Father Flynn - May 27, 54. \\
Bertrams Hartman - to Bertrams; Stolen \\
Clothes Line to Villager - Oct. 55 \\
San Esteban - Times ; \uncertain{Davies} \\
Bedroom fireplace - stolen at Martha Jackson Gallery 1954.}

\marginal{Marg. Bruen. May 54}

\marginal{Esther Baker}

\marginal{Betty + Carl Carmer, '57}

\marginal{another}

\marginal{Stolen by Martha Jackson Gallery 1954.}

\marginal{1954 \quad \uncertain{'53}}

\note{the first three stand above and beside « Jewels of the Madonna. »; the
Martha Jackson matter is written large across the fore-edge margin beside the
Fire Escape and Bertrams Hartman lines, and « another » is interlined in the
Bertrams Hartman line before « Stolen ». The date closing the last line is
written twice, once above and once in the line}

\hand{jo}{\emph{Museum Personelle.} \\
\emph{Boston} \quad Eleanor S. \uncertain{Hunneman} \\
\qquad Assistant to Administration \\
\qquad Edgell \quad - (died) \\
\qquad Constable \\
Robt. Beverley Hale - Met. \\
Frederick Sweet - as. curator - Art Inst. \\
Dr. \uncertain{Lewis} Webster Jones - Pres. Rutgers Univ. \\
Jos? Butler - Butler Arts. Institute, Youngstown, O \\
Thos. Messer - Director of Exhibitions - \emph{Amer. Fed. of Arts} \\
Wadsworth Athenuem - C.C. Cunningham - Director \\
Met. Francis Henry Taylor Dir (died) \\
\emph{Whitney} - Hermon More, Lloyd Goodrich - \\
\qquad Rosalind Irvine, Marg. \uncertain{McKellar}, \\
\qquad John \uncertain{Baurer}. \\
\emph{Detroit} - Edgar P. Richardson - Dir. \\
Columbus Gallery of Fine Arts - Scott B. \uncertain{Malone} - "Dir" \\
Richmond, Va. Mus. of Fine Arts \\
\qquad Leslie Cheek, Jr. married \uncertain{Dopolos Southall} \\
Cochran Mus. Washington \quad Herman Williams \\
Hartford Athenium - \qquad Cunningham}

\marginal{sent post card Xmas. Mar. 51.}

\marginal{Chicago}

\marginal{at Worcester, Mass.}

\marginal{Detroit Inst. of Art}

\marginal{daughter of}

\marginal{Freeman}

\marginal{call \$55}

\note{« Lewis » in the Rutgers line is doubtful stroke by stroke; page 33
writes « Rutgers \quad Lewis W. Jones » plainly, which settles the field of
candidates and not the stroke. The « on » of « Hermon » is written over another
letter, and a double stroke stands under « More ». « Athenuem » and
« Athenium » are both as written, and so is « Cochran » for the Corcoran}

\section{Pages 32--33}

\sheet{96504}{32}
\hand{jo}{Mar. \qquad 54 \quad II \\
Ap. 17, 53 \quad I \\
Rec'd \quad \emph{Reality} -}

\marginal{came out}

\hand{jo}{Appolonio Umbro \quad I \\
Director Dept. / Archives \quad II \\
Biennias de Venezia}

\note{a brace gathers these three lines and carries the two numerals}

\begin{ledgertable}{Y{0.44}Y{0.44}}{ & }
from Jo. Felicia Geffen \quad II & Robt. Beverly Hale \quad I + II \\
Howard Siegh - Msy. \quad II & Van Wyck Brooks \quad II \quad (J) \\
Chas. A. Aitken \quad II & Dr. Stanley Romain Hopper \quad (J) \\
\quad \uncertain{Lu Duble} & A. Donald Adams \quad (J) \quad II \\
Ralph Walker & Henry Seidel Canby \\
E. B. White \quad II & \uncertain{Elmore} D. Clark \quad II \\
Paul \uncertain{Beckley} \quad II & Malcom Cowley \\
Jack Kelley & Robt. Frost \\
Mrs. Walt Kuhn & Marianne \uncertain{Morett} \\
H.C. Bentley, May 14" & Frank Jewett Mather \\
Ruth Emerson \quad II & Ed. Root \quad II \\
Bea Goldsmith \quad II & Stephen Clark \quad II \\
Edgar Kaufmann \quad II & Carl Murchison \\
L. Graves & Earl Ludgin - Chicago \quad II \\
Kruesi's \quad II & Richard Tucker \\
Julian Binford & Philip James \quad II \\
Mrs. Cowles & Guy du Bois \quad II \\
Isabel Whitney & Paul Sachs \quad II \\
Eliz. Taber \quad II & Rutgers \quad Lewis W. Jones \quad I \quad II \\
Clarence Chatterton \quad I II & Jacques Lipchitz \quad Hastings N.Y. \quad II \\
Jno. Biddle \quad II & Glenway Wescott \quad II \\
Edward \uncertain{Nilson} \quad July 53 & \\
Alfred Marx \quad Ag. 53 & Deems Taylor \quad II \\
Cleveland Woodward \quad " & Dr. Richard McCormack \\
\end{ledgertable}

\marginal{+ note}

\marginal{Fredericksbg. Va.}

\marginal{May 28}

\marginal{etc.}

\marginal{etc.}

\note{the two columns are not divided by a rule on this leaf, though page 34's
are; they are given as they stand. Several of the roman numerals are written
over others, and the ticks against Ed. Root, Bea Goldsmith, Kruesi's and Earl
Ludgin are hers. « + note » stands above « Kaufmann », « Fredericksbg. Va. »
after « Binford », « May 28 » above « Whitney », and « etc. » after Canby and
after Clark. A pencil line runs from the left column into the right between the
Aitken and Walker rows, and its two ends cannot be paired}

\section{Pages 34--35}

\sheet{97539}{34}
\hand{jo}{\emph{Reality cont: I} \qquad \emph{II No. only.} \quad II}

\begin{ledgertable}{Y{0.46}Y{0.42}}{ & }
Agnes Tate - June 53 \quad J.H. \quad II & Lloyd Goodrich \\
Amy E. Hogeboom \quad " \quad II & Marie Appleton \quad II \\
Katherine Proctor \quad " & Lewis Mumford \quad III \\
Mr. Cooney at Whitney \quad Nov. 53 \quad J.H. & Jane Grey \quad II \\
Mrs. McCulloch Miller \quad II \quad Oct. 1953 & \struck{Matt} \uncertain{Charles Pearsons} \\
Bertram Hartman \quad II \quad Nov. 18 53. & Rosalind Irvine \quad J.H. + \\
Edwin Bahre \quad II \quad " \quad " & Hugh Ferris \\
Teddy Ballantine \quad II \quad " \quad " \quad " & Norman \uncertain{Herzkel} \\
56 \quad Mahonri Young \quad II \quad " 17" \quad " \quad " & A. Walkowitz \\
Thomas Messer \quad " 26, \quad " \quad " & Dr. Edgell \quad II \\
Mrs. Hobart Nichols - Dec. 5" 53 \quad " & Robt. Beverly Hale \quad I \quad II \\
Louis A. Simons \quad " 4 \quad " & Eliot. Alex. \quad II \\
Marc Connolly \quad II. \quad " \quad " \quad " & \uncertain{Koivumaky} Alex \quad I. II \\
Chas. Hopkinson \quad II \quad " \quad " \quad J.H. \quad " & Lincoln Kirstein \quad II \\
Ruth Underhill \quad Denver \quad " \quad " \quad " \quad " & Mahonri Young Jr. \quad II \\
Francis Henry Taylor \quad Dec. II \quad " 19, 53 & \uncertain{Dmitry Kessarinsky} \\
Elihu Root (Jr) \quad Dec. 19, 53 \quad " & \uncertain{Ada Gabriel} \\
Otto Spaeth \quad " \quad " \quad also " \quad I \quad II & Forbes Watson \quad I - II \\
Russel Lynes - Harpers Mag. \quad " \quad " \quad II & Ida Glackens \\
Hugh McKean - Rollings Coll. Fla. \quad II & Maury Becker \\
28 \quad Walter Pach \quad Jan. 15, 54 & Doris Helm. \\
Edna MacGowan \quad " 17, \quad " & Dorothy Sickler, Art News - \\
Dr. \uncertain{Daghlee Torre} \quad " 19" \quad " & Francis Hackett \quad arts etc. \\
Dr. Slack - Feb \$10 \quad 54. \quad II & Mr. Mrs. Marty Reed \\
Miss Davies - May. 54 \quad II only & Jules + Betty Isaacs \\
 & Mrs. H. B. Arnold Columbus \\
\end{ledgertable}

\marginal{the Grapes-in-law}

\marginal{also \quad Kimball}

\marginal{gallery}

\marginal{arch. - Ant. Kelly \quad lunch}

\marginal{\$5.}

\marginal{Univ.}

\marginal{Ed. H.}

\note{a vertical rule divides the two columns of this leaf, which is why it is
set as a table and page 33's is not. The right-hand column is a second sitting
of ink, smaller, and several of its entries stand between two lines of the left;
the pairing of a right-hand name with the left-hand line beside it is not to be
relied on. « the Grapes-in-law » follows Rosalind Irvine and is struck through
by a long curve; « also » and « Kimball » stand above and below Bahre;
« gallery » above Herzkel; « arch. - Ant. Kelly » and « lunch » above the
Simons line; « \$5. » above Hopkinson; « Univ. » above the Taylor line; and
« Ed. H. » above McKean}

\hand{jo}{9 - x 5 1/8 \emph{large print} - on paper 12 1/4 x 10"}

\marginal{1952.}

\hand{jo}{\emph{Early Sunday A.M. Whitney color print.} \\
Mrs. \qquad Frissell - Mitla - Feb. 53 \\
Mr. Mrs. \uncertain{Mullhrs} \quad Silvia - Truro land - S. Westport - Mass \\
Sam. Cochran, Jr. Bank of N.Y. \quad June 53 \\
Annabell \qquad " \qquad " \qquad 52 \\
\quad Bogdan \quad Colombia Hand laundry \quad 52 \\
- \quad Smith - stole brushes - \quad 52. \\
Gustav Roesch - Halfpenny + Bell - Spring 53. \\
Katherine Proctor \quad 53 \quad Emily Brinker - Oct. 54}

\marginal{Labor Day 52}

\marginal{Summer}

\marginal{Sup.}

\note{« Labor Day 52 » is interlined above « Truro land », « Summer » above
the 52 of the Smith line, and « Sup. » above the 53 of the last, which is
divided from Emily Brinker by a short vertical stroke. A pencil rule closes the
block}

\hand{jo}{II only. \qquad I + II \\
19 \quad Owen Quimby \quad May 54 \\
\quad Harry Sundberg \quad 5, 26, 54 \\
\quad John Dos Passos \\
\quad \uncertain{Reinhold Niebuhr}. \\
\quad \uncertain{Gertrud Kagen}. \\
\quad Bogdan - Laundry \\
\quad Richard \uncertain{Kadman} \qquad Isabel Whitney \$5. Ap. 55. \\
\quad Otto Kallir \qquad \# 10. Henry Allen Moe \\
\quad Bertram Berenson - I + II \qquad Lawrence \uncertain{Cortz} \quad I II \\
\quad Aileen Saarinen \qquad Mrs. Gertrude E. Vogt - Freeman \\
\quad Olive Rush \qquad Mr. Riess - Kennedy Galleries \\
\quad Robertson G. Collins - 164 Oak Drive, Medford \\
\quad \uncertain{Lu Duro}. \qquad (Janet Stark's brother) \\
\quad C.C. Cunningham \qquad Florence Norton, Mrs. Freeman \\
\quad Whitelaw Reid - Tribune \qquad Max Eastman \$10. \\
\quad Otto de Vilbiss \quad II - Wm. Craig. 305 Lexington. \quad Sr.}

\marginal{Harold}

\marginal{4th studio}

\marginal{gallery \quad Fb. Dr. 46 W. 57 St. \quad St. Etienne gal.}

\marginal{1955 May}

\marginal{May 12, 54.}

\marginal{Times}

\marginal{Oregon \quad \$5.}

\marginal{Hartford Athenuem}

\marginal{Danish composer}

\note{« Harold » is interlined above the Quimby line and a second word under it
will not be read; « 4th studio » follows Kadman; the Kallir line's three
interlineations are written over and around the name; « 1955 May » closes the
Moe line and a mark beside it is boxed; « May 12, 54. » stands above Berenson;
« Times » after Saarinen; « Oregon » and « \$5. » above the Collins line;
« Hartford Athenuem » after Cunningham; « Danish composer » above the last
line. This half of the leaf is crowded line against line and the placing of one
or two interlineations against a particular name is not certain}

\section{Pages 36--37}

\sheet{96484}{36}
\hand{jo}{\emph{Reality III} - fall 1955. \\
Bob Stephens - \\
Sydney Weintrob}

\marginal{given to -}

\note{a brace gathers the two names and carries « given to - »}

\hand{jo}{Gift of Bill Johnson - summer 1956. \\
Viva Mexico - Chas. Macomb Flandrau \\
\qquad We read July 55.}

\hand{jo}{\emph{Gift of Katherine Proctor, \uncertain{Charlestown Witt.}} \\
\qquad Sep. 1953. \\
Song of My Life - Yvette Guilbert \\
Voltaire's Genius of Mockery - Victor Thaddeus \\
Arctic Adventure - Peter Freuchen \\
Outline of History - H.G. Wells \\
Life of Michel Angelo - John Addington Symonds \\
Education of Henry Adams - Henry Adams \\
The American Language - H.L. Mencken \\
The \uncertain{Promises} Man Lives By - Harry Scherman \\
The World's Great Letters \quad Ed. by M. Lincoln Schuster \\
Philosopher's Holiday - Irwin Edman \\
Bedside Book of Famous Short Stories \quad Ed. Angus Burrell + Bennet A. Cerf \\
The Ballad + the Source - Rosamond Lehmann. \\
Stories of 3 Decades - Thomas Mann \\
Whale Ships + Whaling - Albert \uncertain{Cork} Church -}

\marginal{Colvin}

\marginal{Enlarged}

\marginal{1936}

\marginal{W.W. Norton Co. 70 - 5 am. \quad 200 photos \quad 1st edition 1938.}

\note{« Colvin » is interlined in the heading between « Katherine » and
« Proctor »; « Enlarged » above « - H.L. Mencken »; « 1936 » above the end of
the Mann line, and the Norton matter is written under a rule across that line
and belongs to the Church book that follows}

\hand{jo}{returned because found inscription inside by Kath. P. \\
" My father, Alonzo Nelson Colvin when a boy of 19 went \\
out from New Bedford as supercargo with Whaler Roscoe. He was \\
gone 5 years having rounded the globe from E + W. to N. + S. \\
Capt. Gorham commanding the Brigue with 30 seamen \quad 1840. \\
Xmas 1938. To be the property of my son Talman C. Budd.}

\note{she opens the quotation and does not close it}

\section{Pages 38--39}

\sheet{96483}{38}
\hand{jo}{\emph{\uncertain{Pintz} straw beads - Mexico 1955} \\
Mrs. Mumm \qquad Rebecca Soyer \\
" \quad Torre \qquad 7 \quad Felicia Geffen \\
Bea G. \qquad Dorothy Becker (Sears) \\
Lu D. \\
5. Niki Pack}

\hand{jo}{\qquad \emph{Both General Elec.} \\
I \quad Elec. blanket - pale - sea green - Altman Jan. 20, 1958. \\
\quad day or 2 before E. return from N.Y. Hos. \\
\quad \# N.P. 153292 - A. \quad Double, \uncertain{sigle} control. \\
II \quad Elec. blanket - darker green \quad Macy - Dec. 1958. \\
\quad not used this winter of 57. 58 \\
\quad \# \\
I = 1957 model \qquad II 1956 \\
I. Catalogue no. BA1 C 22 \quad "Double bed" \\
\quad 115 volts \quad (UL) \quad 190 watts \\
\quad 60 cycles \qquad A.C. only.}

\marginal{foam}

\marginal{control}

\note{« foam » is interlined above « sea green »; the second \# is written and
left without a number. « control » stands at the fore-edge beside the catalogue
line. Small raised marks stand over « Macy », over « Dec. » and over the 9 of
that line's 1958}

\hand{jo}{\emph{More - re Wash. Sq. \# 3.} \\
Hayden Richardson - Hopper studio - \\
\quad back in day of Walter Tittle. \\
He the son of famous architect who built \\
Trinity Church in Boston on Copley Sq. \\
Richardson spread the Romanesque from Coast \\
to Coast - a style copied in all the Carnegie \\
Trinity Boston her windows by LaFarge. \\
Mary Tillinghast - 14 fl. stained glass - \\
\quad worked with LaFarge on windows + by \\
herself also. \\
Richardson Jr. drank. Wife Margaret \\
Lyman Stowe \uncertain{dig} Cor. of Wash. Sq., lived in Judson.}

\marginal{died}

\marginal{Libraries}

\note{« died » is interlined above « Hopper studio » and « Libraries » at the
end of the Carnegie line, which turns without one}

\hand{jo}{Old 10" St. Studio demolished Nov. 1955. \\
Mark Twain House - S.E. cor. 9" + 5" \quad 1954 ± \\
Brevoort + Lafayette \quad Earlier 1953 ? \\
Amer. Engineering Women's House - \\
west side of 5" bet. Wash. Sq. + 8" St. torn down \\
Monsters new apt. ready 1954 \\
Whitney Mus. abandoned 1954 \quad " \\
Rhinlander Houses - Cor. Wash. Sq. + 5. \\
\quad " \quad Terrace - 11" St. S. side W. of 6" Ave. \\
\quad Condemned for new public school 1955. \\
Jewish Cemetery - Triangular plot on W. 11" St. \\
near S.E. Cor. of 6" Ave. Still there winter 1955.}

\note{a long arrow runs up the fore-edge from the Rhinlander line to the
« torn down » that ends the Engineering line}

\section{Pages 40--41}

\sheet{97518}{40}
\hand{jo}{\emph{Antonins Raymond} - DuBois Studio - \uncertain{designed} \\
the lions at sides of fire place. Worked with Frank \\
Lloyd Wright + years in Japan. See New Yorker - \\
met at Architects League ± 1943.}

\marginal{Sept. 22, 56. \quad p. 36, 7.}

\hand{jo}{Albert White Voss - arctic}

\marginal{went to}

\hand{jo}{Elinor \quad --- \qquad Carte Editor. 1" fl. when 13. \\
Sea Craft Real Est. for Blodget Estate \\
Eleanor Wylie \qquad before returned to Sailors Snug \\
\qquad on expiration of lease of land.}

\marginal{Sailors Snug}

\hand{jo}{Negulesco. Jean. rear \\
Schofield Thayer 1" front \\
Frank Harris 2" front \\
John Dos Passos - \\
Andrew Witwer \\
Paul \uncertain{Rosika} \\
\quad --- \quad Lux \\
Gordon + Jane Potter \\
Jos. Pollet \\
E. D. + Louise Moore \\
Margaret Moore (\uncertain{Grunthbach}) \\
Eugene \uncertain{Kierney} \\
Bud \uncertain{Tielman} + \uncertain{Gerdis}}

\hand{jo}{\emph{Sizes} \qquad \emph{Water colors} \\
Purple petunia, white etc. \quad 15 \uncertain{1/2} x 17 1/2 \\
Tree. \quad 20's (22 1/4) x 15 1/2 \\
\quad 1/2 sheet Watmans = 22 1/2 x 15 1/2}

\note{the numerator of the first fraction is a closed loop and reads as a 6}

\hand{jo}{James Mt. (N.M) tree. 18 3/8 x 12 3/4 \\
\qquad gold insert. \\
Ruth Wittenbergs garden \quad 20 x 24. \\
Grand Titons \quad 22 1/2 x 24 + \\
\quad Old mats - outside - 19 x 25 \\
\quad gold insert - inside \quad 14 3/4 x 21. \\
Lady of Good Voyage inside 13 1/2 + x 19 1/2 + \\
Provincetown from hill top \quad 13 1/2 x 20 \\
(smaller) \qquad insert \quad 14 3/4 x 21. \\
Flowers - Peonies \\
\quad White pitcher \\
\quad Black glass \\
Black Eyed Susans in heavy mat 31 3/4 x 26 3/4}

\marginal{1/2 sheet Watman}

\marginal{outside}

\note{a brace gathers the Peonies, White pitcher and Black glass lines and
carries « 1/2 sheet Watman »}

\hand{jo}{\uncertain{Limps} \qquad Vt. trees + mts. \\
\quad 3. \qquad Grand Titon \\
\qquad High tide Gloucester wharf \\
\uncertain{Ponies} + \qquad = 1/2 sheet Watman.}

\note{a brace gathers the three lines against the « 3. »}

\section{Pages 42--43}

\sheet{97162}{42}
\hand{jo}{\emph{Who's Who in American Art} \\
\qquad Vol. \emph{VI} \quad 1956 Edition \\
Mrs. Dorothy B. Gilbert Editor \\
\emph{listing} = \emph{previous} volume.}

\marginal{also 1954.}

\hand{jo}{\emph{Hopper: Jo N. (Mrs. Edward)} - \emph{Painter} \\
3 Washington Square, New York 3, N.Y. \\
B. New York, N.Y. Studied with Robert Henri, \\
Kenneth Hayes Miller. Exhibited: C.G.A. \\
1943; P.A.F.A. 1938; G.G.E. 1939; \\
A.I.C. 1943; B.M.; Weyhe Gal., 1944; \\
Rehn Gal. 1946; C.G.A. 1951; N.A.C. 1950; \\
Church of the Ascension, New York, 1952.}

\note{this block is boxed in red crayon, from above « Hopper: Jo N. » to below
the last line}

\hand{jo}{Published for the Amer. Federation of Arts by the \\
R.R. Bowker Company, 62 West 45 St. \\
New York 36, N.Y. \\
Publishers of Amer. Architects Directory \\
\qquad Amer. Art Directory \\
\qquad Amer. Library \quad " \\
\qquad \emph{Ulrich's Periodicals Directory} \\
this questinaire dated Dec. 15, 1955. \\
\quad Dead line for reply Jan. 3, 1956 \\
Sent to Mrs. Edward Hopper \\
\qquad 3 Wash. Sq. N. \\
\qquad New York 3, N.Y.}

\hand{jo}{this data copied because had to be returned \\
\qquad Jan. 3, 56. \\
\quad Profession Classification - Check list in \\
\quad order of importance. \\
P. (Painter) \\
\qquad Easel - mural - portrait - miniature \\
S. (Sculptor) \quad etc}

\marginal{original}

\note{« original » is interlined above « had to be returned »; a tick stands
after « (Painter) »}

\hand{jo}{Sent invitation for Greenwich Gallery Ap. 58 \\
\quad to Mrs. Gilbert. Note in relation. \\
\quad I to receive next questionnaire when \\
issued. \\
Appeared again in 1958 vol. with additions \\
Greenwich \uncertain{Guladt} Gallery - 1958, 1959.}

\marginal{Jan.}

\marginal{water colors etc.}

\marginal{(into one man show)}

\section{Pages 44--45}

\sheet{97471}{44}
\hand{jo}{\emph{Amer. Artist} - Jan. 1956 \\
contains - Portrait by E. on Coal box near press - \\
\qquad by Sidney Weintraub - Budd Studio \\
Cape Cod A.M. in color - page. \\
10 A.M. \quad --- \quad Dawn in Pa. - \quad 2 on the Aisle \\
McComb's Dam Bridge - 11 A.M. - Room in N.Y. \\
House by R.R. - \qquad \emph{Sunday, Hoboken}}

\marginal{E. H.}

\marginal{full}

\note{« E. H. » stands above « Portrait » and « full » above « page. »}

\hand{jo}{\qquad Sent to \\
S. Cochran. Feb. 17, 56 \\
E.G. Robinson \qquad " 22 \\
Dr. A.M. McLellan \quad " 24 \\
Mr. Mrs. Robt. Slater \quad " \quad "}

\note{a vertical stroke closes these four lines at the right}

\hand{jo}{\qquad In previous Buick, 1939 - acquired 1941.}

\entry{1941 -}{1941}
\hand{jo}{Las Vegas, Los Angeles, Frisco, Portland, O. - Cape}

\marginal{Columbia R.}

\entry{1943 -}{1943}
\hand{jo}{Mexico - City, Saltillo, Sept. Monterrey}

\marginal{by R.R.!}

\marginal{no Cape}

\entry{1946}{1946}
\hand{jo}{Mexico - Sattillo \quad Grand Titon \quad Cape.}

\marginal{Buick 39}

\entry{1942 -}{1942}
\hand{jo}{Early - J. hemeroid operation at St. Clare Hos. \\
\qquad Dr. Felix Sullivan}

\note{a brace gathers the 1943, 1946 and 1942 lines at the left. « Saltillo »
and « Sattillo » are both as written}

\entry{1948.}{1948-01}
\hand{jo}{Jan. Group of 6 or 7 paintings at Rehn. \\
\quad very successful.}

\marginal{E. H.}

\entry{1948}{1948-02}
\hand{jo}{Feb. N.Y. Hospital - Prostate \\
\qquad Grosvenor Hotel when left Hospital \\
\qquad Stayed into May \\
3 trips to N.Y. Hos. hemmorhages. \quad Cape}

\marginal{Dr. Allister McLellan}

\marginal{On this occasion}

\entry{1949 -}{1949}
\hand{jo}{Back to N.Y. H. - 2" stage of prostate \\
operations. Dr. McLellan. \quad Cape.}

\entry{1950}{1950}
\hand{jo}{N.Y. Hospital - Diverticulosis \\
\quad Dr. \uncertain{Zugerman} got Dr. John G. Schmidt. \\
Feb. - Retrospective Ex. Hopper, Man of the Year \\
\quad at Whitney Mus. - W. 8" St. \quad E. in hospital \\
at opening of Exhib. \\
\quad Exhib. moved to Boston Mus. of Fine Arts + \\
3 day visits at Col. Rich. Magee on Bunker Hill \\
Exhib. on to Detroit - on return from \\
Chicago Art Institute - E. degree D.F.A. \\
stayed at Palmer House. + parties. \\
\quad Cape + after Detroit trip - night on train}

\entry{1951 -}{1951}
\hand{jo}{Mexico - Sattillo in heat wave. On to \\
Santa Fé + return to Cape - began efforts for \\
Electric. \quad Cape 1952.}

\entry{Dec. 1952}{1952-12}
\hand{jo}{Water Color Show at Metrop. J.H. head of Bertrams Hartman \\
\qquad Sum. 1952. \\
\qquad Dec. 7 or 6 off for Mexico - via \\
El Paso, TX. + South. Stay at \uncertain{Guanacinato} (Posada \\
de la Presa, Mitla in Oaxaca - S. of \ill{} at}

\marginal{Next day off to Mexico}

\note{the last line runs into the fore-edge and its close cannot be read. The
1948 to 1952 memoranda run forwards, unlike the chronicle of pages 19 to 24}

\keywords{frames and framing, picture sizes, water colour paper, Frank K. M.
Rehn, Heydenryk, The Studio Gallery, American Artists Group monographs,
Broadway theatres, insurance policies, Mexico, South Truro, Washington Square,
Reality, Oliver Baker, photographs of works, Who's Who in American Art,
museum staff, Whitney Museum of American Art, Greenwich Village demolitions,
Josephine Hopper's own work, gift books}
\end{document}
